\documentclass[../main.tex]{subfiles}
\begin{document}
%==========================================TIÊU ĐỀ TẬP========================================
\begin{table}[h]
    \begin{adjustwidth}{-1cm}{}
        \begin{tabular}{>{\centering\arraybackslash}p{6cm}|>{\raggedright\arraybackslash}p{12.5cm}}
            \multirow{5}{6cm}
            % Cột bên trái: ÉPISODE và số tập
            {\\[-40pt]
                \flaregothic\fontsize{20pt}{20pt}\selectfont \centering
                \color{eptype}HISTOIRE\color{black}\\
                \freshbot\fontsize{72pt}{72pt}\selectfont
                \color{epnum}1
                \\[18pt]
                \flaregothic\fontsize{14pt}{14pt}\selectfont
                \color{epattr}BIENVENUE, \uline{\color{Mano} Aloe} !
                \color{black}
            }
             &
            % Dòng thứ nhất: tên tập tiếng Nhật
            {\fotskip\fontsize{18pt}{18pt}\selectfont \ruby{家}{いえ}に\ruby{潜}{ひそ}むサキュバス}
            \\[6pt]
            % Dòng thứ hai: tên tập tiếng Anh
             & {\cafeteria\fontsize{18pt}{18pt}\selectfont Succubus in da House!?}            \\[4pt]
            % Dòng thứ ba: tên tập tiếng Việt
             & {\vnmsans\fontsize{18pt}{18pt}\selectfont Chuyện về nàng Succubus trong nhà}   \\[8pt]
            % Dòng thứ tư: Nhân vật xuất hiện
             & {\brian{\fontsize{14pt}{14pt}\selectfont Track used: Dark Ages - Brain Busters}
               }
            \\[8pt]
            % Dòng cuối cùng: Thời gian diễn ra của tập (thường sẽ là ngày phát hành)
             & {\continuum\fontsize{16pt}{16pt}\selectfont Ngày phát hành: 07/02/2021}        \\[18pt]
        \end{tabular}
    \end{adjustwidth}
\end{table}
%=========================================TRUYỆN CHÍNH========================================
\color{black}

\begin{center}
    ({\brian Music used: Dark Ages - Ultimate Battle})
\end{center}

\begin{center}
    \uline{(Tiếp tục từ {\flaregothic ÉPISODE 92}, quyển số Bảy.)\\Chiều tối hôm đó.}
\end{center}

\pov{Hikawa Kuon}

Trời đổ mưa tầm tã, những hạt mưa nặng trịch rơi lộp bộp trên mặt đường bê tông, loang loáng dưới ánh đèn đường mờ mờ. Không khí ẩm ướt quấn lấy người tôi, khiến từng hơi thở cũng mang theo chút se lạnh. Tôi vừa chạy vừa kéo cái cặp máy tính chống nước lên che đầu, dù biết rằng chẳng mấy tác dụng, cũng sợ rằng cái máy tính của tôi sẽ lại dính nước. May thay, trạm xe buýt không quá xa công ty, nếu không thì có lẽ giờ này tôi đã ướt như chuột lội mất rồi.

Đứng dưới mái hiên của trạm xe, tôi phủi nhẹ mấy giọt nước mưa bám trên áo khoác, rồi khẽ thở dài. Công việc cuối cùng của năm đã hoàn tất, bao gồm cả việc phân bua cho cả Botan-chan và Luna-chan, đồng thời giải quyết việc viện phí cho Marine-chan. Giờ đây, tôi chỉ còn việc đón chuyến xe này để về nhà ăn Tết.

Tôi mở điện thoại, định đọc báo giết thời gian trong lúc chờ chuyến xe đến, thì màn hình bỗng nhấp nháy một cuộc gọi. Là ông sếp, Tanigo-san. Chắc lại gọi để hỏi thăm với chúc Tết tôi đây mà.

Tôi nhấc máy. ``Dạ? Hikawa Kuon xin nghe?''

Giọng ông sếp vang lên bên kia đầu dây, vẫn trầm trầm như mọi khi. ``Hikawa-san đấy à? Hôm nay cậu định về nhà ăn Tết sao?''

Tôi hơi nhíu mày. ``Dạ\dots\ vâng? Hôm qua cháu báo cho chú biết rồi mà?''

``Về sớm thế!''

\dots Không biết ông ấy có đang đùa không, nhưng tôi chỉ im lặng chờ ông nói tiếp. Đã qua ngày Ông Công Ông Táo rồi, còn sớm sủa gì nữa.

``Tôi đang định tìm cậu để báo một tin này. Nhưng trước tiên\dots''

``Vâng?''

``Tóm tắt cho tôi nghe nhóm \textit{holofive} đến công ty làm việc thế nào.''

Tôi ngả người ra sau, tựa vào cột trạm xe buýt, mắt nhìn xa xăm qua làn mưa. Những thành viên thế hệ Năm à? Cũng lâu rồi tôi chưa tổng hợp lại nhỉ.

Xem nào\dots

``Polka-chan, Lamy-chan và Botan-chan có màn 'chào hàng' khá là\dots\ nói thế nào nhỉ, quái dị: Pol-pol tự nhận mình là mèo (hoặc đúng hơn, bị Fubu-chan xúi giục), Lamy-chan định biến văn phòng thành tủ đựng rượu, còn Botan-chan thì đạp xe ba bánh như một đứa con nít.''

Nói đến đây, tôi khẽ cười nhẹ. Chỉ ba câu thôi mà đã đủ tóm gọn sự hỗn loạn của những ngày hôm ấy trong ba tuần vừa rồi.

Nàng Cáo sa mạc và tôi thì chơi trò giải cứu con tin với Fubuki-chan, Ayame-chan, Haato-chan với Okayu-chan và Koro-chan làm những nhân vật chính.

Nàng Yêu tinh tuyết thì được tôi rủ về nhà uống rượu với bố tôi và đã có một cuộc tỉ thí rượu.

Nàng Sư tử thì\dots\ làm trò cười trước mặt Towa-chan và Luna-chan.

Sau khi tôi nói xong, bên kia đầu dây bỗng im lặng một hồi lâu. Tôi hơi nghiêng đầu.

``Vậy là còn hai người cậu chưa gặp nhỉ,'' Tanigo-san lên tiếng.

Tôi thở nhẹ. ``Thực ra là một ạ. Một người nữa thì bị chính chú cho nghỉ chỉ sau khoảng hai tuần stream từ tháng Tám năm ngoái rồi. Hoặc ít nhất là cháu nghe vậy.''

``Ừ. Tiếc cho cô ấy thật. Giá mà cô ấy cẩn thận hơn thì\dots'' giọng ông sếp thoáng chút tiếc nuối, nhưng rồi cũng không nói tiếp. Tôi gật đầu, mặc dù biết ông không thể thấy.

``Mà thôi, chuyện đã qua rồi ạ. Người còn lại cháu nghĩ là Momo\dots\ ờm, gì đó, Nene. Botan-chan bảo mùa xuân tới mới lên văn phòng làm việc ạ.''

``Rồi, biết thế.'' Ông đáp. ``À mà, đấy, đó là điều mà tôi muốn báo cho cậu này.''

Tôi nhướng mày.

``Vâng\dots? Chú cứ nói tiếp.''

Tôi có cảm tưởng rằng những điều mà Tanigo-san sắp nói ra đây sẽ là những điều rất quan trọng. Tôi ghé sát điện thoại hơn vào tai.

Bầu không khí bỗng chốc căng thẳng đến lạ.

``Về người mà cậu nói là bị cho nghỉ việc ấy\dots\ Thực ra là không hẳn là cô ấy bị đuổi đâu.''

Vậy là tôi đoán sai à? Nhưng dù gì thì, đó vẫn chưa phải là những gì ông ấy muốn nói hết.

``Tôi sẽ nói cho cậu biết điều này, nhưng cậu phải hứa với tôi là phải tuyệt đối giữ bí mật, được không?''

Có gì đó không ổn ở đây. Có việc gì mà ông ấy đang giấu kín tới mức như thể là một cái gì đó tuyệt mật vậy?

``\dots Kể cả là bố mẹ cháu ạ?'' tôi nhướng mày hỏi lại.

``Miễn là cậu bảo họ không được nói cho ai biết là được. Tôi cũng đã nói chuyện này cho Yuujin A-san rồi, và cô ấy cũng đã giữ kín miệng, không cho idol \hll\ biết.''

Thông tin này có lẽ là chỉ những người thuộc tầng quản lý cấp cao mới được biết thôi. Nó thực sự có một điều gì đó tối quan trọng, như thể sẽ quyết định vận mệnh của cả công ty này.

Tôi hít một hơi thật sâu, chuẩn bị nhận tin tuyệt mật từ CEO của chúng tôi. ``\dots Đã rõ. Chú cứ nói ạ.''

\ct

\begin{center}
    Đường \ruby{Kaihou}{Giải Phóng}, quận \ruby{Choushimai}{Hai Bà Trưng}, Kawachi.
    Trên xe buýt, tuyến J7\footnote{Trong câu chuyện này, tuyến J7 chạy cùng đường với tuyến 26 ở Hà Nội (Mai Động - SVĐ Quốc Gia); tuyến đi qua Tokyo Dome - Thị trấn Điện tử (nơi mà văn phòng \hll\ tọa lạc) - đường Kaihou - Bến xe Bitei (Mỹ Đình ở ngoài đời).}.
\end{center}

``Điểm dừng tiếp theo: Đại học Bách khoa Kawachi.''

Tôi tựa đầu vào cửa kính, mắt nhìn những giọt nước mưa lăn dài trên mặt kính như những sợi tơ mỏng. Ngoài trời vẫn mưa không ngớt, những ánh đèn đường phản chiếu lên mặt đường ướt át, kéo dài thành những vệt sáng mờ ảo.

Tôi thở ra một hơi thật chậm, mắt vẫn nhìn ra ngoài mà chẳng thực sự để tâm đến cảnh vật.

Đúng như tôi dự đoán. Ông sếp Tanigo-san của chúng tôi lại bày ra một điều bất ngờ ngay trước thềm năm mới. Nhưng lần này, điều đó xảy ra ngay tại chính nhà tôi.

Tôi không phải người dễ bị bất ngờ. Làm trong \hll\ một thời gian đủ dài, tôi đã quá quen với những màn ``bày trò'' của công ty - từ việc quản lý lịch trình của các thần tượng cho đến những quyết định đột ngột từ cấp trên. Nhưng lần này, mọi chuyện lại diễn ra một cách quá bí ẩn.

Tôi biết rõ rằng toàn bộ \hll\ đang tập luyện cho \textit{Bloom,} - buổi live concert lớn sắp tới. Với lịch trình dày đặc như thế, gần như không có cơ hội nào để ai trong số họ xuất hiện ở nhà tôi. Vậy thì\dots\ người đang chờ tôi ở đó là ai?

Câu trả lời rõ ràng nhất chính là người mà công ty đã cho nghỉ việc.

Dựa trên những gì mà ông ấy nói, người đó không hẳn bị đuổi, mà thực tế là tự xin nghỉ, rồi được công ty chuyển sang một công việc khác. Nhưng nếu đã như vậy, tại sao em ấy lại muốn gặp tôi? Và còn ngay tại nhà tôi?

Bố mẹ tôi có biết chuyện này không? Nếu biết, tại sao họ không nói gì với tôi? Hay là Tanigo-san đã nhờ họ giữ bí mật?

Tôi khẽ nhíu mày, tự hỏi liệu ông sếp già của mình còn đang toan tính điều gì nữa mà tôi chưa biết.

Mà thôi\dots\ chuyện đã đến nước này thì tôi cũng chẳng thể làm gì ngoài việc chờ đợi.

Tôi tiếp tục ngồi lặng trên xe buýt, đầu óc cứ xoay vòng trong những câu hỏi không lời giải. Tôi đắm chìm trong suy nghĩ đến mức suýt nữa thì lỡ mất điểm dừng.

Chỉ đến khi loa thông báo trạm tiếp theo, tôi mới giật mình, vội vàng đứng dậy, chuẩn bị xuống xe.

\begin{center}
    \uline{15 phút sau.\\Nhà Hikawa - quận Kyuukami, Kawachi.}
\end{center}

Gió lạnh quất vào mặt tôi ngay khi bước xuống xe buýt, mưa phùn cứ rả rích mãi không thôi. Ánh đèn đường vàng vọt chẳng đủ sưởi ấm gì ngoài việc soi rõ những vệt nước long lanh đọng lại trên mặt đường ẩm ướt. Tôi rùng mình, kéo chặt áo khoác, vừa run vì lạnh vừa lẩm bẩm trong đầu: về nhà vào cái thời tiết này đúng là thảm họa.

Không ô, không mũ, quần áo thì ướt sũng, cảm giác chẳng khác gì bị ném thẳng vào một cái tủ đông. Răng tôi va vào nhau lập cập khi đặt chân đến cửa nhà. Cuối cùng cũng về được\dots\ chuẩn bị cho một mùa Tết không kèn không trống nữa-

``\vie{Kuon ơi! Lại đây xách phụ bố với bác đống đồ đạc này cái!}'' tiếng của bố tôi cất lên ngay trước cửa nhà.

À phải, tôi còn chưa kịp nghỉ ngơi, mà phải giải quyết chuyện này trước đã.

Trước mặt tôi là một đống đồ lỉnh kỉnh, có vẻ như hai người khiêng còn không xuể. Tôi đảo mắt một lượt qua mấy cái hộp, rồi dừng lại khi thấy một cái hộp có nhãn ``phụ kiện máy tính để bàn.''

Khoan\dots\ nhà mình mua máy tính mới cho ai à? Tôi chưa từng nghe bố mẹ nói gì về việc này. Nếu đó là cho tôi thì cũng không cần thiết, tôi xài máy tính xách tay quen rồi, chẳng muốn đổi làm gì.

Nghĩ vậy, tôi đang định mở miệng nói thì bố tôi tiếp lời: ``\vie{Xách lên tầng hai con nhé!}''

Tầng hai? Tôi hơi cau mày.

``\vie{Cái phòng trống đấy ạ? Nhưng bên nào ạ?}'' tôi hỏi lại.

``\vie{Cầu thang bên trái ấy.}'' Bên nhà tôi sao?

Tôi khựng lại một chút. Khoan đã\dots

Tôi chưa nói với các bạn nhỉ? Cả hai phía tầng hai nhà tôi đều đang bỏ không từ sau đại dịch năm ngoái. Trước đó, tầng ấy từng được một cửa hàng viễn thông thuê, nhưng sau cùng cũng phải trả lại mặt bằng. Bố và bác tôi thậm chí còn đang định cho người khác thuê lại\dots

Vậy mà giờ, lại có người chuyển vào?

``\vie{Mày xách cẩn thận nha, kẻo lại hỏng đồ của con bé đấy!}'' một giọng người đàn ông vang lên, hớt hải cùng bố tôi bê đồ.

Con bé? Tôi hơi ngớ người.

``\vie{Trong này có gì hả bác?}'' tôi để cặp máy tính lên một chiếc thùng và bắt đầu bốc lên, hỏi.

Bác ấy đáp lại: ``\vie{Toàn cốc với bát đũa cả đấy! Để cặp máy tính của mày lên đấy rồi xách một thể!}''

Tiện thể tôi cũng giới thiệu luôn: người mà xưng hô ``tao-mày'' với tôi bằng giọng hơi khàn ấy là Hikawa Hijaki (\ruby{火}{ひ}\ruby{川}{かわ}\ruby{日}{ひ}\ruby{邪}{じゃ}\ruby{気}{き}, \ruby{Hi}{Hỏa}-\ruby{ca-oa}{Xuyên}\ \ruby{Hi}{Nhật}-\ruby{gia}{Da}-\ruby{ki}{Khí}) - anh trai của bố tôi. Người sống cùng nhà với chúng tôi ở tầng dưới, cùng với hai người con gái hiện đang làm xa nhà, và bà nội tôi. Bác ấy là con trai thứ của nhà Hikawa - trên bác còn một người anh cả (đã qua đời cách đó khoảng ba tháng trước, đúng ngày mà tôi phải bỏ dở đám tang chỉ để sửa lại Robo-PON\footnote{Xem lại {\flaregothic ÉPISODE 84}, quyển số Bảy.}), nhưng nếu hỏi về ông ấy thì câu trả lời thường là: ``Thôi đừng nhắc, tính xấu nhiều quá đến mức không muốn nói tới.''

Bác Hijaki-san xưa nay vẫn quen xưng ``mày - tao'' với tất cả mọi người, tôi cũng quen với điều đó rồi.

Nhưng có một chuyện quan trọng hơn. ``\vie{Nhà mình có người thuê mới ạ?}''

``\vie{Ừ. Mà bố nghĩ chắc là Tanigo-san nói với con rồi chứ?}''

Điều đó có nghĩa là ông ấy cũng đã nói chuyện với gia đình tôi rồi sao? Tôi vừa cảm thấy bối rối, lại vừa nghi ngờ, nhất là với tình cảnh hiện tại bây giờ.

``\vie{Ông ấy vừa gọi điện cho con ban nãy xong, nhưng lại chỉ nói rằng đó là một người nào đó đang đợi ở nhà\dots}'' tôi đáp lại.

Bố tôi gật gù, hất đầu lên cầu thang: ``\vie{Bố nghĩ chắc là em ấy đấy.}''

\dots Vậy là đúng rồi.

Cái người mà Tanigo-san bảo tôi giữ bí mật tuyệt đối, giờ đang ở ngay trong nhà tôi.

``\vie{Thôi thôi, nhanh nhanh lên, kẻo hỏng đồ của người ta bây giờ!}'' bác tôi lên tiếng, ngay trước khi tôi lại rơi vào trạng thái trầm tư.

Tôi không có thời gian để suy nghĩ thêm. Chúng tôi hối hả khiêng đống đồ lên gác, bác tôi đi trước, rồi đến bố, cuối cùng là tôi. Với ba người, thế là vừa đủ để vận chuyển hết đống hành lý đó.

Nhưng dù chân tôi vẫn bước, đầu tôi thì lại cứ xoay vòng trong những câu hỏi không có lời giải.

\ct

Không mất quá nhiều thời gian để chúng tôi mang hết đồ của vị khách thuê mới lên căn phòng chỉ định. Mấy cái thùng không quá nặng, nhưng kích thước cồng kềnh khiến việc di chuyển có chút vất vả. Tôi vừa đặt một thùng xuống sàn thì bỗng nghe thấy giọng một người con gái.

``Các bác để đấy đi, để cháu tự xếp lại cho. Mấy người già rồi, đụng vào lỡ làm bể thì sao\~{}?'' Một giọng điệu chảnh chọe, nhưng cũng pha chút tinh nghịch.

Bác Hijaki-san đáp lại, thở dài: ``Mày xếp được không đấy? Tao thấy phòng đây chật lắm.''

``Hứ, bác đang coi thường cháu hả? Chẳng phải chuyện của bác đâu, để đấy cho cháu!'' cô gái ấy đáp lại.

Vừa nói, cô gái vừa khoanh tay, hất cằm đầy tự tin. Tôi nhướng mày ngạc nhiên với cách nói chuyện có phần\dots\ hỗn xược này. Y như Shion-chan hồi tôi mới vào làm.

``Mà nè, anh kia tới chưa vậy bác?'' cô gái ấy lém lỉnh hỏi.

Bố tôi nhíu mày, nhìn về phía tôi: ``Cháu đùa à? Nó đang lên đây này. Chuẩn bị có hàng xóm mới rồi đấy.''

``Gì chứ\~{}? Tất cả cũng tại ông YAGOO chứ ai! Đáng lẽ cháu phải được ở chung với anh ấy luôn mới phải! Còn để cháu qua đây làm gì không biết nữa!''

\dots Sao cứ như thể em ấy bị ép qua đây vậy?

Tôi cẩn thận đặt thùng cuối cùng xuống, rồi hướng về phía bố tôi. ``Ờm\dots\ Con để đống này ở đâu vậy bố?''

Nhưng thay vì bố tôi trả lời, cô gái kỳ lạ kia đáp hộ bằng một tông giọng bề trên: ``Ở kia! Nhanh lên coi, đứng chắn đường người ta mãi vậy hả?''

Tôi nhíu mày, rồi thở dài làm theo: ``Ờ ờ, biết rồi, khổ quá cơ\dots\ Bình tĩnh nào\dots''

Tôi nhẹ nhàng đặt thùng xốp đồ sứ xuống góc phòng, rồi mới quay người lại. Và ngay khoảnh khắc đó, tôi sững sờ.

Đứng trước mặt tôi là một cô gái với mái tóc dài màu hồng cùng với chút tiêu điểm màu vàng ở bên trong, cặp mắt xanh lục sắc sảo trông đầy tinh quái. Trên đầu em ấy là một đôi sừng không cân đối, phía sau còn có một chiếc đuôi dài đang quấn hờ lấy một cái micro.

Chất giọng kênh kiệu, cái cách đứng khoanh tay đầy tự tin, cùng vẻ mặt tinh nghịch. Một succubus, không nghi ngờ gì nữa.

Giống Choco-san thật đấy. Ngoại trừ việc em ấy là một con quỷ hấp tinh thực thụ.

Tôi chưa kịp mở miệng thì em ấy đã lên tiếng trước: ``Nhìn gì em ghê thế\~{}? Anh chắc là Hikawa Kuon-kun nhỉ?''

\dots Xưng hô tôi với ``-kun'' ngay từ lần đầu gặp mặt? Hiếm thấy thật. Cô gái này có lẽ không phải là dạng bình thường, không chỉ theo nghĩa về ngoại hình, mà là cả về tính cách.

``Ờm\dots\ Ừ. Còn em là\dots''

Cô gái nhếch môi cười. ``Thử đoán coi nào\~{}? Hay là ngay cả mặt em mà cũng không nhận ra hả?''

Tôi khẽ nheo mắt. Khuôn mặt này\dots\ quen lắm. Tôi chắc chắn đã từng thấy ở đâu đó trước đây.

Sơ yếu lý lịch của em ấy\dots?

Tôi bắt đầu hỏi. ``Em cũng là Quỷ như Towa-chan và Choco-san à?''

``Rõ rành rành thế này mà còn phải hỏi à? Anh bị cận thị hay gì?''

Tôi nhún vai đáp lại: ``Hỏi cho chắc thôi,'' và hắng giọng một tiếng. ``Thế hệ thứ Năm đúng không?''

Cô nàng nhướng mày, khoanh tay lại.

``Đáng lẽ là thế\dots\ Nếu như ta không lỡ chân tự đạp vào số phận ngay buổi stream đầu tiên.'' Giọng điệu nghe vừa thản nhiên, vừa có chút cay cú.

``Hừm. Kỹ năng đặc biệt của em là gì?''

Cô nàng lập tức giơ tay lên đầy tự hào: ``Tất nhiên là hát rồi! Ta đâu phải loại succubus rảnh rỗi chạy qua thế giới loài người chỉ để chơi!'' em ấy hất mặt tự đắc về phía chúng tôi. ``Ta sẽ trở thành một diva lẫy lừng cả Quỷ giới!''

Tôi khẽ nheo mắt, nhìn cô nàng với vẻ nghi ngờ. ``\dots Ờm, OK.''

Lập tức, cô nàng Succubus ấy nổi quạu: ``Cái thái độ đó là sao hả? Anh dám xem thường em à!?''

Tôi giơ hai tay làm bộ đầu hàng: ``Không có gì. Anh chỉ cần hỏi em mấy câu như thế thôi, hình như anh nhớ ra rồi.''

Tôi chậm rãi suy ngẫm lại các dữ kiện.

Một succubus. Có ước mơ trở thành nữ danh ca. Cùng thế hệ với Polka-chan, Botan-chan, Lamy-chan và một thành viên nữa. Xưng hô kiểu ``anh-em'' ngay từ lần đầu gặp một cách đầy kiêu ngạo.

Chỉ có duy nhất một người khớp với tất cả những điều đó.

Tôi hít sâu một hơi rồi nhìn thẳng vào cô gái trước mặt.

``\dots Mano Aloe-chan (\ruby{魔}{ま}\ruby{乃}{の}アロエ, Ma-\ruby{nô}{Nãi/Ái}\ A-lô-e\footnote{Còn có thể đọc là A-rô-e.}), đúng chứ?''

Cô nàng lập tức chống tay vào hông, vênh mặt lên: ``Pín-pon! Chuẩn không cần chỉnh luôn!'' Bố tôi và bác tôi nhìn nhau, rồi đồng loạt quay sang tôi ngạc nhiên:

\begin{dialogue}
    \speak{Hijaki} \vie{Mày từng gặp nó ở đâu rồi à?}
    \speak{Kuon} \vie{Dạ không, cháu chỉ dùng suy luận thôi ạ.}
    \speak{Ryujin} \vie{Suy luận kiểu gì, thử giải thích xem.}
\end{dialogue}

Thực ra, suy luận ra cái này không khó.

Tôi bình tĩnh tóm tắt lại: \textit{holofive}, hay Thế hệ thứ Năm, lẽ ra có năm thành viên. Trong ba tuần vừa rồi, tôi đã gặp ba người, người còn lại dự kiến sẽ đến vào mùa xuân tới.

Bằng phương pháp loại trừ, cộng với các manh mối từ cách nói chuyện của cô nàng trước mặt, tôi có thể đoán được danh tính của em ấy: Mano Aloe. Một Succubus không cùng họ hàng với Choco-san, ranh ma hơn cả Towa-chan, và tuyên bố muốn trở thành một diva nơi em ấy sống (hoặc ít nhất là em ấy nói vậy).

\chrint

\begin{wrapfigure}[28]{r}{8.5cm}
    \includegraphics[width=8.5cm]{../../../../Image/Book?/aloe_model.png}
\end{wrapfigure}

Giờ tôi mới nhớ ra trong bản CV mà tôi đã từng đọc qua, Aloe-chan là một cô nàng Succubus đang được hướng luyện với tính cách bướng bỉnh, tinh nghịch và có thái độ bề trên. Một người dễ dàng chỉ trích người khác, nhưng lại không chịu được khi người khác chỉ trích ngược lại em ấy.

Ước mơ của cô nàng này không gì hơn là trở thành một Diva ở xứ Quỷ giới - nơi quê nhà của em ấy. Hoặc ít nhất, là đã từng là như vậy.

Aloe-chan có cách nói chuyện theo kiểu của một đứa nhóc được chiều hư hỏng - hay ``kusogaki'', giống như Shion-chan - thậm chí có phần huênh hoang và thất lễ với bề trên. Điều này thì tôi không ngạc nhiên, vì dường như bất cứ ai ở xứ Quỷ cũng đều ngạo mạn như thế. Đến cả Ayame-chan lúc mới lên lần đầu cũng gọi người xem của em ấy là ``đám con người'' chứ không hề nể nang gì.

Nhưng ngoài những tính xấu ra, em ấy là một người thích nấu ăn - dù tôi không chắc là em ấy có nấu ăn giỏi không, hay tôi sẽ lại phải đối mặt với một Haachama phiên bản 2.0. Ngoài ra, em ấy là một người thích dọn dẹp nhà cửa\dots\ nhà người khác, chứ tuyệt nhiên lại không phòng ổ của chính mình.

Là một nữ quỷ, nên những bộ phim kinh dị - nhất là ở trên P*ime - là sở trường của em ấy. Nhưng những trò chơi cùng loại thì\dots\ trớ trêu thay, lại không như vậy. Em ấy cũng là một người ghét sâu bọ, kể cả loài bướm, nhưng em ấy khẳng định sẽ không giết chúng. Thật kỳ quái.

Một điều kỳ lạ nữa, là dù có giọng hát thiên bẩm - tự nhận, chúng tôi chưa kiểm chứng điều đó - nhưng Aloe-chan lại thích sử dụng các công cụ điều chỉnh giọng nói như autotune và giọng ``kero kero'' để bắt chước giọng tổng hợp như Vocaloid khi hát.

Em ấy có một cặp sừng quỷ không cân xứng (với bên trái dài và cong hơn bên phải). Em ấy khẳng định rằng sừng như vậy vì em ấy vẫn đang phát triển, thậm chí khăng khăng nói rằng đến một lúc nào đó cả hai chiếc sừng sẽ có cùng kích thước. Ngoài ra, Aloe-chan còn có một cái đuôi quỷ dài, cuốn lấy một cái míc-rô. Quanh cổ em ấy có một cái tai nghe trùm đầu. Em để tóc xõa, dài đến giữa lưng, có màu hồng bên ngoài, bên trong là màu vàng, trên mái có chiếc cài hình những ngôi sao bốn cánh màu tím, trên sừng trái còn có đeo nơ tím-hường.

Ngay lúc này đây, Aloe-chan diện một chiếc áo sơ mi trắng không tay có cổ, đeo một cái thắt ngực được buộc dọc với áo, trước ngực là một cái ruy-băng màu hồng gắn cùng với một vật trang trí hình xương; mặc một bộ váy đầm xếp nếp màu tím bên ngoài và màu hồng bên trong, cùng với những sợi tưa lưa màu đen bên trong. Ở bên hông váy phải có đồ trang trí hình trái tim màu hồng. Aloe-chan đi hai chiếc tất đen (bên trái) và tím có ca-rô trắng (bên phải), đi giày cao gót đen có mũi và đế màu hồng tím.

À mà, tôi chưa nói rằng em ấy có mặc một chiếc quần lót đen bên trong chứ nhỉ? Mà, tại sao tôi lại nói đến cái đó?

E hèm.

\chrint

Tôi vẫn cảm thấy có cái gì đó hơi cấn cấn\dots\ Không phải vì việc Aloe-chan là một succubus hay cái thái độ kiêu ngạo, mà là vì chuyện em ấy lại xuất hiện ngay tại đây, trong nhà của chúng tôi, và rồi trở thành hàng xóm - thậm chí là người nhà - của gia đình chúng tôi. Tại sao YAGOO-san lại sắp xếp như vậy chứ?

``\vie{Ông ấy có nói gì với bố chuyện này không?}'' tôi quay đầu về bố tôi hỏi.

Bố tôi lắc đầu: ``\vie{Chịu. Hôm qua ông ta gọi cho bố, bảo là sẽ có người từ công ty con chuyển vào ở chung, nhưng lại chẳng nói gì thêm cả.}''

Tôi cau mày, tỏ vẻ nghi ngờ với ý định của ông ấy. Càng lúc tôi càng có nhiều câu hỏi hơn cần được giải đáp.

Câu hỏi đầu tiên, và cũng là quan trọng nhất mà tôi muốn được biết là: ``\vie{\dots Tại sao em ấy lại dọn sang nhà mình sống?}''

Tôi hỏi vậy với hai người đàn ông lớn tuổi hơn, nhưng rồi nhận lại chỉ là cái nhún vai và cái lắc đầu. ``\vie{Bó tay. Mà cái đó thì sao không hỏi thẳng em ấy luôn đi? Hỏi bố, bố sao biết được?}''

\dots Cũng phải. Tốt hơn hết là cứ hỏi người đang có mặt ở đây.

Dù vậy, tôi vẫn cảm thấy hơi lấn cấn. Tôi tiến về phía cô nàng succubus vẫn đang chống hông đầy tự mãn kia. ``Tại sao em lại chuyển về đây sống cùng với bọn anh?''

Aloe-chan nhướng mày, giọng đầy khiêu khích: ``Anh cần gì phải biết? Liên quan gì đến anh à?''

Tôi đổ mồ hôi hột với cái thái độ lồi lõm đó. Dù sao thì, tôi cũng đã biết (hoặc chí ít, vừa mới nhớ ra) rằng cô nàng này đúng kiểu ``một đứa ranh'' chính hiệu.

Và với một đứa ranh như thế, thì sẽ cần có một cách tiếp cận phù hợp hơn.

``Đương nhiên là liên quan chứ\dots\ Không thì anh hỏi làm cái cụ gì?''

Nàng Succubus thấy vậy, đáp trả: ``Được thôi. Anh muốn nghe thì cũng được. Dỏng tai lên mà nghe em nói cho rõ đây!''

Tôi nhún vai: ``Anh không bị điếc đâu mà sợ. Nói đi.''

``Với cả để ý cách ăn nói chút đi. Cháu như vừa rồi là hơi hỗn đấy,'' bố tôi lên tiếng nhắc nhở.

Thấy vậy, tôi lầm bầm, tỏ vẻ cam chịu như đã quá quen với điều này: ``Đến Towa-chan cũng không hách dịch đến thế này đâu. Mà thôi, kệ em ấy đi bố. Quỷ là thế mà.''

Aloe-chan hắng giọng, khoanh tay lại đầy tự mãn và bắt đầu giải thích - hoặc có lẽ là em ấy cố gắng như vậy. ``Mọi người cũng đã biết ta muốn trở thành một diva của Quỷ giới rồi, đúng chứ?''

Cả ba chúng tôi gật đầu.

``Và vì ta là một succubus chính hiệu, nên ta sẽ-''

Tôi lập tức cắt lời em ấy: ``Sử dụng giọng hát của chính mình để mê hoặc mọi người, theo cả nghĩa tinh thần lẫn thể xác chứ gì?''

Aloe-chan ngay lập tức nổi cáu: ``\textbf{Để em nói hết đi nào! Làm tụt hết cảm xúc của người ta!}''

Tôi bật cười: ``Xin lỗi, xin lỗi. Thói quen cướp cò của anh hồi học cấp ba thôi.''

Em ấy lườm tôi, phồng má lên vì bị tôi cướp lời: ``Mà anh nghĩ cái quái gì thế hả! Mê hoặc về tâm hồn thì đúng, chứ thể xác là cái quỷ gì chứ!?''

Tôi giả vờ ngơ ngác đáp lại: ``Anh tưởng succubus nào chả thế? Loài succubus như em toàn chuyên thu hút cánh đàn ông và\dots\ ờm, làm gì đó với họ mà, em biết chứ?''

Lập tức, mặt em ấy đỏ hỏn lên, rồi nhảy dựng như bị tôi bắt bài: ``\textbf{Em không phải là loại succubus kiểu như vậy!}''

Tôi nhìn Aloe-chan chằm chằm, rồi cười khẽ: ``Nhìn mặt đỏ hỏn thế kia là biết rồi. Anh nói chuẩn chưa?''

Em ấy lắp bắp, gương mặt vẫn đỏ ửng lên: ``K-Không có! Đừng có giỡn mặt với em nữa!''

Phản ứng ngượng ngùng của em ấy dù là một succubus càng làm tôi cười lớn hơn. Tất nhiên là bác Hijaki-san và bố tôi thì không thoải mái gì khi tôi bất lịch sự như vậy.

``Thôi, trêu thế đủ rồi. Nói tiếp đi.''

Mano-chan trấn tĩnh lại, hắng giọng rồi tiếp tục với giọng đầy tự mãn: ``Ta rất tự tin với khả năng này, nên ta đã đến thế giới loài người, để cho các người phải rung động và phục tùng trước ta!''

Em ấy kết thúc bài phát biểu bằng một nụ cười đầy kiêu hãnh.

\dots Nhưng tất cả những gì em ấy nhận lại chỉ là ánh mắt đầy nghi hoặc của ba người chúng tôi.

Bác Hijaki-san khẽ nhướng mày, khoanh tay lại cất lời: ``Đó mới chỉ là mục đích mày đến đây, thế nhưng tại sao mày lại chọn nơi này để sống cùng thì lại chưa nói.''

Tôi tiếp lời, tay chống cằm suy nghĩ: ``Thật. Thà rằng tìm một cái phòng trọ khác còn hơn chỗ chật chội này\dots'' rồi nghiêng đầu nhìn em ấy. ``Chắc phải liên quan gì đó đến YAGOO-san chứ, đúng không?''

Aloe-chan trố mắt, nhìn tôi như thể vừa nghe thấy chuyện gì đó động trời: ``S-Sao anh biết!?''

Tôi không đáp ngay mà chỉ nhướng mày, rồi trả lời bằng một câu hỏi khác: ``Ủa, ông ấy chưa nói với em rằng anh là tổng quản lý của \hll\ à?''

Cô nàng Succubus chớp mắt mấy lần, rồi hét lên: ``\textbf{Cái này ta có biết đâu!}''

Tôi suýt nữa thì phun cả nước ra. ``Bảo sao trông em bất ngờ thế. Mà chắc ông ta lại đùn công việc cho mình đây mà.''

Aloe-chan vẫn còn ngơ ngác, rồi vội vã giải thích: ``Ta chỉ biết rằng ông ấy cho ta địa chỉ ``ký túc xá \hll\'', và nó ghi trên bản đồ là ở đây thôi!

Cả nhà Hikawa chúng tôi đồng loạt sững người, rồi cùng nhau lặp lại.

``Ký túc xá á?''

Tôi nhíu mày nhìn bố tôi và bác tôi: ``\vie{Sao con lại không biết nhỉ?}''

Bố tôi lắc đầu, khuôn mặt vẫn bày tỏ sự ngạc nhiên: ``\vie{Bố cũng thấy bất ngờ đấy.}''

Bác tôi thì giơ tay lên như thể xin hàng: ``\vie{Cái này thì tao chịu. Việc đấy là của nhà chúng mày. Tao không liên quan gì hết nha.}''

Tôi lắc đầu, ngửa mặt lên trần nhà: ``\vie{Chắc là ông ta lại bày trò gì đây mà\dots\ Tí nữa con hỏi lại ông ấy.}''

``\vie{Ừ. Bố cũng thấy kỳ lạ\dots}''

Tôi ngước nhìn em ấy lại một lần nữa. Khuôn mặt em ấy rõ ràng em ấy đang không nói dối.

Càng lúc mọi thứ càng khó hiểu và rối rắm hơn rồi đấy\dots

\ct

\begin{center}
    [{\color{red}\textbf{CẢNH BÁO:}} Đoạn sau đây đề cập đến một sự kiện nhạy cảm liên quan đến sự nghiệp của Mano Aloe. Nếu bạn cảm thấy không thoải mái với chủ đề này, hãy cân nhắc trước khi tiếp tục đọc.]\\({\brian Music used: Dark Ages - Cerebrawl, but without the kazoo part})
\end{center}

Về việc ``ký túc xá'' này, bố con tôi sẽ tìm hiểu sau. Cả ba người đàn ông chúng tôi, cùng với Aloe-chan, tiếp tục ngồi nghe kể chuyện. Nhưng bây giờ, ngoài việc tại sao em ấy lại ở đây, tôi có một thắc mắc khác, mà có lẽ là sẽ hơi\dots\ khó nói.

``Cái chuyện ký túc xá túc xủng này thì tạm gác đã. Bây giờ anh có một cái thắc mắc: tại sao Tanigo-san lại đưa em vào đây vào lúc này, chứ không phải từ nửa năm trước?''

Aloe-chan hơi giật mình khi nghe câu hỏi ấy. Biểu cảm tinh ranh và tự mãn ban nãy vụt tắt, thay vào đó là một ánh mắt chùng xuống, phảng phất sự đắn đo.

``\dots Cái này khó nói lắm. Em cũng\dots\ không biết nên bắt đầu từ đâu nữa.''

Tông giọng của em ấy trầm hẳn đi, không còn chút vui tươi hay nghịch ngợm như trước. Giờ đây, nó phảng phất sự tiếc nuối, một chút buồn bã, và cả sự hối hận.

``Anh biết là em đã xin nghỉ từ tháng Tám năm ngoái. Nhưng em đã đi đâu và làm gì trong sáu tháng sau đó?'' tôi hỏi, dù lần này có phần dè dặt hơn.

Quả thật đó không phải là câu hỏi mà tôi nên đặt vào lúc này.

Có vẻ như bố tôi và bác Hijaki-san cũng có chung những thắc mắc với tôi:

\begin{dialogue}
    \speak{Ryujin} Tại sao cháu nghỉ làm vậy?
    \speak{Hijaki} Có cái gì cứ kể hết nghe đi xem nào. Biết đâu nó có thể giúp ích một chút.
\end{dialogue}

Aloe khẽ gật đầu, ngẩng mặt lên, đôi môi khẽ gợn lên một nụ cười méo xệch: ``\dots Cảm ơn ba người.''

Em ấy hít một hơi thật sâu, rồi bắt đầu trải lòng. ``Nói thật, ở trên thế gian này, ta cũng chỉ tầm tuổi một học sinh trung học thôi. Ta như kiểu\dots\ như lũ nhân gian các người hay nói gì nhỉ, `tấm chiếu chưa trải', đúng không?''

Cả ba chúng tôi gật gù. Quả nhiên, nếu tính theo tuổi người thì em ấy cũng chỉ có mười lăm tuổi, ít tuổi hơn cả em gái tôi.

``Tao cũng nghĩ mọi người ở đây nói như vậy.''

``Chả trách mà chú thấy cháu trẻ thế.''

``Đó. Cũng giống như thanh thiếu niên các người, ta cũng muốn được tỏa sáng giống như nhân gian các người nữa!'' em ấy tiếp tục. ``Ta biết rất rõ nền idol ở thế giới này đây, biết rõ mọi người muốn gì, và\dots''

Trước khi để em ấy nói tiếp, tôi lên tiếng: ``Cái này anh phải công nhận đúng. Nền idol ở đây đã phát triển đến mức tất cả mọi người, không kể lứa tuổi hay gốc gác, cũng muốn cống hiến cho nó, đúng chứ?''

Aloe-chan gật đầu. Thật sự với một người mới chỉ từng này tuổi mà đã có ước mơ lớn lao như thế thì đó là một điều hiếm khó có tìm.

``Và với bộ âm thanh mà bọn chú mang lên ban nãy'' - chúng tôi hướng mắt về đống hộp ở góc phòng mà chúng tôi đã khiêng lên ban nãy, rồi lại nhìn em ấy - ``chắc là\dots\ cháu cũng thích hát lắm nhỉ?''

``Đúng vậy. Ta cũng biết năng lực mình ở đâu, và vì vậy, ta đã cố gắng hết sức, hết sức có thể, để có thể được gia nhập vào \hll.''

``Nhìn qua thì chắc mày cũng đã cố hết sức để có thể được vào công ty đó, đúng chứ?''

``Dạ vâng\dots\ Sau bao ngày vất vả, cuối cùng ta cũng đã được tuyển vào đó, cùng với Polka, Nene, Lamy và Botan.'' Em ấy bắt đầu nói hết những gì trong tâm tư của mình một cách chân thành nhất, thậm chí còn nói chuyện như là những người ``cùng đẳng cấp'' chứ không còn hống hách như vừa nãy.

Em ấy nở một nụ cười yếu ớt. Một nụ cười chất chứa đầy những kỷ niệm, hy vọng, và cả niềm kiêu hãnh về một thành tựu mà mình đã đạt được. Nhưng chỉ trong thoáng chốc, nụ cười ấy nhạt dần.

``Ta đã nghĩ rằng ta có thể bắt đầu cống hiến cả tài năng của mình cho mọi người, đã có thể bước chân vào sự nghiệp, nhưng rồi\dots''

Chợt đến đây, em ấy ngừng lại, sụt sịt, mắt bắt đầu đẫm lệ.

``\dots một tai nạn đã xảy đến với ta.''

Cả ba người đàn ông chúng tôi chững lại. Không gian trong phòng bỗng trở nên yên lặng đến đáng sợ.

``Hôm đó\dots\ ta đã phạm phải một sai lầm nghiêm trọng.''

Em ấy nói với giọng run rẩy. Tôi thấy hai bàn tay em ấy khẽ siết chặt lại, như thể đang cố kiềm chế một cơn xúc động mạnh.

``Ta đã làm một buổi stream thử nghiệm, chỉ đơn giản là để kiểm tra âm thanh và các hiệu ứng\dots\ nhưng ta không biết rằng nó đã bị bật công khai thay vì ở chế độ riêng tư.''

Em ấy cười khổ, ánh mắt xa xăm như đang hồi tưởng lại giây phút định mệnh ấy. Một sai sót nhỏ, nhưng đã phải đánh đổi lại một thứ gì đó quá khủng khiếp.

``Không chỉ có thế\dots\ trong lúc đó, ta lại còn lỡ nói một số điều mà đáng lẽ ta không nên nói. Về công ty, về dự án, và\dots\ cả những chuyện riêng tư của bản thân.''

Bác Hijaki-san và bố tôi liếc nhìn nhau, như thể đã phần nào đoán ra được vấn đề, mặc dù họ có kiến thức khá mơ hồ về nền giải trí.

``Lúc đó ta không nghĩ gì nhiều. Chỉ đơn giản là đang thử nghiệm thôi mà! Nhưng ta đã quên mất một điều quan trọng: Internet không bao giờ tha thứ cho bất kỳ sai lầm nào.''

Giọng em ấy nghẹn lại.

``Chỉ trong vòng vài giờ, đoạn clip ấy bị lan truyền khắp nơi. Nó bị cắt xén, bị trích dẫn, bị bóp méo. Ta trở thành chủ đề bàn tán trên khắp các diễn đàn. Và rồi\dots''

Em ấy ngừng lại, cố nuốt xuống nghẹn ngào.

``\dots mọi thứ bắt đầu vượt khỏi tầm kiểm soát.''

Em ấy rùng mình khi nhớ lại.

``Ta bị chỉ trích, bị đe dọa, bị tấn công tinh thần. Những tin nhắn, những lời lẽ cay nghiệt\dots\ nó dồn dập đến mức ta không còn biết phải làm gì nữa.''

Mắt em ấy bắt đầu đỏ hoe.

``Họ lục lọi quá khứ của ta. Họ vẽ ra những câu chuyện mà ta chưa từng làm. Họ nguyền rủa ta, hạ nhục ta, và thậm chí còn\dots''

Giọng em ấy lạc đi.

``Ta đã phải khóa tài khoản, phải tránh xa mạng xã hội\dots\ Nhưng ngay cả như vậy, nó cũng không dừng lại.''

Cả ba người chúng tôi im lặng, không nói một lời nào trước những lời bộc bạch đáng sợ này. Tôi nắm chặt tay lại. Những gì em ấy đã trải qua\dots\ tôi không thể tưởng tượng nổi nó kinh khủng đến mức nào.

``Và cuối cùng\dots\ ta đã không chịu nổi nữa. Ta\dots\ rời bỏ mọi thứ.''

Không gian chìm vào tĩnh lặng.

``Ta cứ nghĩ rằng, nếu ta biến mất, nếu ta rời khỏi \hll, thì mọi chuyện sẽ kết thúc. Nhưng\dots''

Một giọt nước mắt lăn dài trên má em ấy.

``\dots nó không dễ dàng như vậy.''

Tôi nhìn em ấy, trong lòng dâng lên một cảm giác khó tả. Không còn là một con succubus kiêu ngạo và nghịch ngợm nữa. Trước mặt tôi bây giờ chỉ là một cô gái trẻ, đã từng ôm ấp ước mơ, nhưng lại bị vùi dập bởi những thứ mà em ấy chưa từng lường trước.

Và tôi biết, cuộc hành trình của em ấy vẫn chưa kết thúc.

``Ngay từ lúc đó\dots\ ta đã biết, ta đã hiểu rằng\dots\ mình đã không còn đường lui nữa.''

Aloe-chan bật cười khổ, nhưng trong ánh mắt lại chỉ toàn là nỗi đau.

``Ban đầu ta còn cố nghĩ rằng\dots\ `Ừ, chắc vài ngày nữa mọi chuyện sẽ lắng xuống thôi.' Nhưng rồi, đó lại là một suy nghĩ cực kỳ sai lầm. Càng ngày, mọi thứ càng trở nên tồi tệ hơn.''
Giọng em ấy trở nên run rẩy, đôi bàn tay siết chặt vào vạt áo của chính mình.

``Người ta không chỉ dừng lại ở việc công kích ta trên mạng. Họ bắt đầu tìm kiếm thông tin về ta, tìm mọi cách để moi ra quá khứ của ta, và rồi\dots\ họ đã tìm được.''

Một khoảng lặng nặng nề bao trùm căn phòng.

``\dots Trong đó có cả bạn trai cũ của ta.''

Cả ba người chúng tôi cau mày. Tôi biết, ở trong giới giải trí idol thì không có chỗ cho sự yêu đương, nhưng\dots\ đó lại là điều mà tôi không ngờ tới.

Tôi định cất lời hỏi, nhưng\dots\ tôi cảm tưởng rằng đây không phải là thời điểm thích hợp làm việc này. Tôi chỉ im lặng, tiếp tục ngồi em ấy kể lại sự tình.

``Hắn\dots\ đã lợi dụng cơ hội này để đổ thêm dầu vào lửa.''

Em ấy cắn môi, mắt cụp xuống như thể không muốn đối diện với ký ức ấy nữa.

``Hắn đã tung ra những chuyện riêng tư giữa ta và hắn, đã xuyên tạc những gì ta từng nói, từng làm\dots\ như thể ta là một kẻ đáng bị hủy hoại. Hắn\dots\ hắn biết rõ ta không thể phản kháng.''

Cơn giận bùng lên trong tôi, nhưng tôi kiềm chế lại. Giờ không phải lúc.

Nhưng có vẻ như, bác trai thứ của gia đình tôi thì lại không thể nhẫn nhịn được. Ngay khi em ấy dứt câu, bác ấy buông ra một câu: ``Đúng là một thằng khốn nạn\dots''

Tôi vội đứng lên, trấn an tinh thần bác: ``\vie{Thôi, thôi bác ơi. Bình tĩnh, bình tĩnh ạ. Cháu cũng đang sôi máu đây, nhưng mà bác để em ấy kể hết đã ạ.}''

Bố tôi thì nghiêng đầu, tỏ vẻ không tin vào những gì mà ông ấy nghe: ``Đến mức đó cơ à?''

Aloe-chan gật đầu, nước mắt long tròng kể tiếp: ``Không chỉ có thế\dots\ Những người ghét ta đã lợi dụng điều đó để tấn công mạnh hơn nữa. Họ nói ta không xứng đáng làm idol. Họ bảo ta là đồ giả tạo. Họ nói\dots\ rằng ta đáng bị như thế.''

Mỗi chữ em ấy thốt ra đều như những nhát dao cứa vào tim mình.

``Và rồi, ta bắt đầu sợ hãi. Sợ mở điện thoại. Sợ nghe tiếng thông báo. Sợ cả việc ra khỏi nh\dots\ Vì ta biết, bên ngoài kia, luôn có những người muốn nhìn thấy ta gục ngã.''

Em ấy siết chặt tay, cố kìm lại dòng cảm xúc đang cuộn trào.

``Ta không ăn được, không ngủ được. Ta chỉ biết khóc mỗi ngày. Cứ mỗi lần ta nghĩ đến việc quay lại, nghĩ đến việc bước tiếp\dots\ thì hàng trăm, hàng ngàn lời cay độc lại ập tới.''

Một giọt nước mắt rơi xuống. Rồi hai giọt, ba giọt, và rồi chẳng mấy chốc, một hàng lệ đã rơi xuống. Nước mắt từng lúc nhiều hơn, và em ấy không thể kìm nén nổi nữa.

``Ta cảm thấy mình vô dụng. Ta tự trách bản thân. Tại sao ta lại ngu ngốc đến vậy? Tại sao ta lại để mọi chuyện đi đến mức này? Tại sao\dots''

Aloe-chan nghẹn ngào, giọng nói dần lạc đi vì cảm xúc hỗn độn trong lòng.

``\dots tại sao ta lại tồn tại? Ta tồn tại\dots\ để làm gì?''

Lồng ngực tôi thắt lại. Hai người lớn trong gia đình Hikawa chúng tôi cũng trở nên nghiêm túc hơn bao giờ hết.

``Ta đã nghĩ đến việc\dots\ chấm dứt tất cả.''

Tôi cứng người lại. Một cảm giác lạnh tới sống lưng.

``Nhưng rồi, ta nhớ đến một thứ\dots\ Ta nhớ đến những người đã luôn ủng hộ ta. Ta nhớ đến những giấc mơ của mình. Ta nhớ rằng, ta vẫn muốn cất giọng hát, vẫn muốn đứng trên sân khấu.''

Em ấy lau vội nước mắt, nhưng chẳng thể ngăn được dòng lệ tuôn rơi.

``Nhưng ta không thể làm thế trong hoàn cảnh đó. Ta đã quá kiệt sức. Ta chỉ muốn trốn khỏi tất cả. Và rồi\dots\ ta quyết định rời đi\dots''

Em ấy không nói hết câu, cảm xúc của em ấy dâng trào, Aloe bật khóc nức nở và sà vào lòng với người gần nhất với em ấy - trớ trêu thay, lại là tôi.

Tôi cũng chỉ có thể dỗ dành em ấy bình tâm lại, nhẹ nhàng nói: ``Mọi chuyện đã qua rồi, em cũng đã chịu khổ quá nhiều, nhất là với một người còn có nhiều dự định phía trước\dots''

Cả bố và bác tôi vừa an ủi, tỏ vẻ tiếc nuối cho em ấy, vừa tức giận trước hành động vô nhân tính của kẻ tự xưng là bạn trai em ấy.

Cả căn phòng im bặt hẳn lại, chỉ còn tiếng khóc ai oán của vị khách mới đến này.

Tôi khẽ thở dài, siết nhẹ bờ vai em ấy như một sự trấn an. Câu chuyện mà em ấy vừa kể không khác gì một cơn ác mộng giữa đời thực, một cơn ác mộng không ai đáng phải trải qua. Tôi liếc nhìn bố và bác tôi, cả hai đều có vẻ mặt nặng trĩu.

Tôi khẽ cười, không phải vì vui vẻ, mà là vì cay đắng cho số phận tréo ngoe của em ấy.

``Anh thật sự không thể tưởng tượng nổi em đã phải chịu đựng những gì. Một mình em chống chọi với cả thế giới, với những con người độc ác ngoài kia\dots\ Đó là một điều quá sức đối với một người trẻ như em.''

Aloe-chan vẫn vùi mặt vào ngực tôi, bờ vai nhỏ nhắn của em ấy vẫn run lên.

``Nhưng em biết không? Em vẫn ở đây. Em vẫn tồn tại. Em đã chọn tiếp tục sống, bất chấp tất cả những gì đã xảy ra. Điều đó\dots\ đủ để chứng minh rằng em mạnh mẽ hơn bất cứ ai ngoài kia.''

Bác Hijaki-san hừ một tiếng đầy phẫn uất, khoanh tay lại và nói: ``Tao nói thật, nếu là tao, chắc tao đã đấm bỏ mẹ cái thằng khốn nạn đó rồi. Còn đám người đã từng hành hạ mày nữa, đệch mợ đúng là lũ rác rưởi mà.''

Bố tôi lắc đầu, ánh mắt đượm buồn nhưng giọng nói lại rất chân thành: ``Chú không thể thay đổi quá khứ của cháu, nhưng chú tin rằng cháu vẫn còn tương lai phía trước.''

Ông ấy liếc về bộ dàn âm thanh mà em ấy mang tới lần nữa, rồi nhìn em ấy. ``Cháu đã từng muốn đứng trên sân khấu, đã từng muốn cống hiến hết mình, đúng không? Vậy thì\dots\ đừng từ bỏ giấc mơ của mình.''

Tôi vỗ nhẹ lưng của Aloe-chan. ``Đúng vậy. Không ai có thể thay đổi quá khứ, nhưng hiện tại và tương lai vẫn nằm trong tay em. Em không nợ ai bất cứ điều gì cả. Em chỉ cần sống vì chính mình, vì những gì em yêu thích mà thôi.''

Nàng Succubus ngẩng mặt lên, đôi mắt đỏ hoe nhìn chúng tôi. ``\dots Thật sao?''

Tôi gật đầu: ``Ừ. Em không đơn độc. Có những người vẫn tin tưởng và chờ đợi em. Và bây giờ, bọn anh cũng ở đây.''

Bờ môi em ấy khẽ run run. Một giọt nước mắt nữa lại lăn dài trên gò má, nhưng lần này, tôi có cảm giác rằng đó không chỉ là nỗi buồn, mà còn có cả sự nhẹ nhõm. Một sự nhẹ nhõm vì cuối cùng đến bây giờ, em ấy có được những người đồng hành.

``\dots Cảm ơn\dots\ mọi người\dots''

Tôi khẽ cười, trong lòng bỗng thấy thanh thản hơn một chút.

Cả căn phòng vẫn im lặng, nhưng sự nặng nề khi nãy đã vơi đi phần nào. Tôi biết, vết thương trong lòng Aloe-chan sẽ không thể lành ngay lập tức, nhưng ít nhất\dots\ em ấy đã không còn phải chịu đựng một mình nữa.

\begin{center}
    \uline{[Kết thúc phân đoạn nhạy cảm.]\\{(\brian Kazoo enabled)}}
\end{center}

\ct

Sau một hồi để Aloe-chan giải tỏa tinh thần của mình, đến mức mà áo của tôi ướt đẫm lệ hòa cùng với nước mưa ngoài trời, em ấy cuối cùng cũng bình tĩnh lại. Em ngồi đó, hai tay ôm cốc nước ấm mà tôi vừa đưa cho, khẽ nhấp một ngụm. Khuôn mặt em ấy vẫn còn lấm lem nước mắt, đôi mắt đỏ hoe, nhưng ánh nhìn đã không còn trống rỗng như trước.

``Hẳn là em đã phải trải qua nhiều khó khăn rồi,'' tôi nói, tay rót cốc nước ấm.

Bác Hijaki-san lắc đầu, nhấp một ngụm trà, giọng trầm xuống như đang suy ngẫm điều gì đó: ``\vie{Giới trẻ bây giờ đúng là tệ hại thật.}''

Tôi thở dài, chống cằm, ánh mắt hướng về phía nàng Succubus, rồi nói chậm rãi: ``\vie{Xã hội là thế, bác ạ. Đôi khi chỉ vì lợi ích bản thân mà có kẻ sẵn sàng đạp đổ nhân phẩm của người khác mà không hề chớp mắt.}''

``\vie{Đúng là đau đớn thật. Khổ thân con bé\dots}'' bố tôi thở dài, rít thêm một điếu thuốc lá.

Aloe-chan siết chặt cốc nước trong tay, giọng khàn khàn cất lên: ``Chứ còn sao nữa! Trong suốt ba tháng sau đó, ta đã bị trầm cảm đến mức tuyệt vọng. Có những lần ta đã muốn tự tử rồi cơ\dots''

Tôi lập tức nín lặng, tim như thót lại. Tôi có thể cảm nhận được bầu không khí trong phòng lập tức đông cứng, trở nên nặng nề đến mức khiến người ta nghẹt thở.

``Yeesh.''

Nhưng chưa dừng lại ở đó, em ấy tiếp tục nói, giọng càng lúc càng nhỏ: ``\dots Và ta đã có một lần suýt thành công như thế\dots''

Cả ba người chúng tôi giật mình, ánh mắt tôi mở to nhìn em ấy, nhưng vị khách của chúng tôi có vẻ không còn quan tâm nữa.

``Cháu\dots\ nói sao cơ!?'' / ``Mày nói thật đấy à!?'' / ``Thật sao!?'' cả ba chúng tôi đồng thanh.

Aloe-chan nở một nụ cười méo mó, khẽ rũ mắt xuống: ``Nếu như hôm đó bạn cùng trọ với ta không phát hiện ra, có lẽ ta đã bị ngạt thở mà lên gặp thiên thần rồi đấy. Hay nói vui hơn là, gặp Kanata-senpai rồi.''

Bố tôi và bác Hijaki lập tức lặng người, nhìn nhau với vẻ mặt khó tin. Tôi cũng không cần hỏi nhiều hơn nữa, chỉ cần nghe đến đó là tôi đã hiểu chuyện gì đã xảy ra. Tôi nuốt khan, không biết phải tiếp tục cuộc trò chuyện này như thế nào.

``Cho chú hỏi\dots\ Cháu đã vượt qua khoảng thời gian khó khăn đó thế nào vậy?'' bố tôi ngập ngừng lên tiếng, dường như muốn chuyển chủ đề.

Nàng Succubus chớp mắt một cái, khẽ nhíu mày như thể đang hồi tưởng lại: ``Không giấu gì mọi người, ta cũng chỉ mới bình ổn được tâm lí cách đây hai tháng thôi. Kể từ khi ta sống ở đây, Kawachi. Ban đầu ta cũng sợ sệt lắm, môi trường ở đây khác hẳn so với thành phố Điện tử, chứ chưa nói là xứ Quỷ nữa. Nhưng rồi\dots''

``Nhưng sao?'' đến lượt bác Hijaki-san hỏi. Một sự căng thằng kéo dài đến ngạt thở, sợ rằng em ấy sẽ lại gặp điều gì đó khó khăn.

Nhưng không. Aloe-chan khẽ mỉm cười, lần đầu tiên kể từ lúc em ấy bắt đầu câu chuyện này: ``Mọi người ở đây khác hẳn so với hồi ở thành phố Điện tử, họ tốt bụng lắm!''

Tôi hơi nhướng mày, rồi cười nhẹ: ``Đặc sản quê anh mà. Họ không quan tâm em là ai hay đến từ nơi nào đâu, họ vẫn sẽ chào đón em một cách đầy nhiệt tình đấy thôi.''

``Cũng nhờ có họ chịu lắng nghe những điều khó khăn mà ta đang trải qua, chẳng bao lâu sau ta mới có thể bình thường nói chuyện với mọi người như thế này rồi đấy!'' em ấy tiếp tục, khuôn mặt dần tươi tỉnh.

Bố tôi khẽ thở phào nhẹ nhõm, nét mặt ông đã bớt đi phần nào căng thẳng: ``May quá rồi.''

Nàng Succubus đặt cốc nước xuống, ngả người ra sau, khoanh tay lại, hớn hở: ``Với cả nghe tin thằng cha ấy bị bắt vào tù giam, ta cũng hả dạ lắm.''

Đến lượt tôi và bác Hijaki-san cũng giãn cơ mặt ra. Tên đó cuối cùng cũng đã có được một kết cục đầy xứng đáng. ``Kết thúc có hậu rồi ha. Vượt qua được sóng gió như vậy là ô-kê-la rồi đấy.''

Tôi nhìn gương mặt vui vẻ của em ấy, tiếp tục giở giọng: ``Cơ mà\dots''

``Sao? Gì nữa?''

Tôi nghiêng đầu nhìn em ấy, nheo mắt lại: ``\dots Trông em vẫn còn lanh chanh lắm.''

Aloe-chan lập tức gườm gườm nhìn tôi, nói với tông giọng của bề trên như trước: ``Anh có ý kiến ý cò gì hả? Em cho anh cái đạp bây giờ\dots''

Bố tôi lắc đầu bật cười vỗ nhẹ vào vai em ấy: ``\vie{Nào Kuon, đừng trêu em.}''

Tôi cũng khẽ cười khi thấy Aloe bắt đầu trở lại với bản tính hay đốp chát của mình. Mặc dù câu chuyện của em vừa kể vẫn còn ám ảnh trong tâm trí tôi, nhưng ít nhất bây giờ, em ấy đã có thể cười nói như trước. Đó là một điều tốt.

``Không có ý kiến gì đâu. Chỉ là\dots\ anh mừng vì em có thể nói chuyện như thế này thôi.''

Bác Hijaki-san nhấp một ngụm trà, ánh mắt ông trầm ngâm như đang suy xét con người của cô nàng này: ``Mày đúng là mạnh mẽ thật đấy. Không phải ai cũng có thể bước qua chuyện như thế mà tiếp tục sống. Nhất lại là con gái.''

Aloe hừ nhẹ, khoanh tay lại. ``Đương nhiên rồi! Ta đây là ai chứ!? Một nữ quỷ hùng mạnh đến từ Dị giới cơ mà! Một cú vấp ngã như vậy thì làm sao mà quật ngã được ta chứ!?''

``Thế mà vừa nãy em còn cho bọn anh thấy con người yếu đuối của em ấy,'' tôi bông đùa. Ngay sau đó, em ấy lườm tôi, mặt tối sầm lại, trông như thể sẽ ăn tươi nuốt sống tôi.

``Quên đi những gì anh thấy đi.''

Bố tôi khẽ lắc đầu, nhưng đôi mắt vẫn ánh lên vẻ ấm áp: ``Cháu đúng là một đứa trẻ kiên cường. Nhưng nếu có chuyện gì\dots\ cháu không cần phải một mình gánh chịu nữa đâu.''

Aloe-chan khẽ giật mình, đôi mắt mở to nhìn ông.

``Đúng thế. Bây giờ em không còn một mình nữa. Dù thế nào đi chăng nữa, vẫn có người sẵn sàng giúp em, lắng nghe em. Và đó không chỉ là bọn anh đâu, mà có khi là mọi người ở thành phố Kawachi này đây!''

Aloe chớp mắt vài lần, rồi bất chợt ngoảnh mặt đi, như thể không muốn ai thấy đôi mắt em ấy đang dần đỏ hoe.

``\dots Các người\dots đúng là phiền phức.''

Chúng tôi bật cưởi với phản ứng có phần tsundere của em ấy.

\begin{dialogue}
    \speak{Hijaki} Biết sao được, bọn tao lớn tuổi rồi mà.
    \speak{Kuon} Thôi nào, anh biết tỏng là em đang vui mừng ra mặt kìa.
    \speak{Aloe} I-Im đi.
\end{dialogue}

Cả căn phòng trầm mặc trong giây lát, nhưng không còn bầu không khí nặng nề như trước nữa. Những cơn mưa bên ngoài cửa sổ đã dần ngớt, nhường chỗ cho màn đêm yên tĩnh.

\ct

Sau một hồi nói về quá khứ đau buồn, chúng tôi quyết định đổi chủ đề để xoa dịu không khí nặng nề. Nàng Succubus tóc hồng cũng lấy lại tinh thần nhanh chóng, ánh mắt dần lóe lên vẻ tinh ranh vốn có. Em ấy hắng giọng một cái, rồi tiếp tục kể chuyện với giọng điệu chảnh chọe như lúc đầu, như thể chưa từng có gì xảy ra.

Bác Hijaki-san khoanh tay, chậm rãi hỏi: ``Vậy sau khi mày rời công ty thì mày đã làm gì?''

Aloe-chan khoanh tay, chống cằm, tỏ vẻ ta đây rất bất mãn: ``Ta cũng đã dự định quay trở lại Quỷ giới sau khi rời \hll\ đấy, nhưng nếu thế thì\dots nhục nhã lắm! Ta không chấp nhận như vậy!''

Tôi khẽ nhướng mày, dự cảm được phần nào câu chuyện tiếp theo. Dù sao thì, nhìn cách em ấy hăng hái kể lể, có lẽ em ấy đã phải lăn lộn không ít để sống sót trong những ngày ải khổ đau ấy.

Tôi mỉm cười nhẹ: ``Và để anh đoán, em quyết định bám víu nơi đây để kiếm sống với niềm đam mê của em, đúng không?''

Aloe-chan búng ngón tay cái ``chóc'' một cái, cười rạng rỡ: ``Không sai! Từ khi ta từ biệt với bọn Gen 5, ta đã phải ở trọ nhiều nơi ở Thị trấn Điện tử, làm đủ nghề để trang trải cuộc sống.''

Bố tôi chống tay lên đầu gối, nhìn cô gái nhỏ bé trước mặt với vẻ tò mò: ``Vậy cháu làm những nghề gì?''

`Ban đầu ta cũng chỉ làm những công việc vụn vặt như quét dọn với lao công thôi. Nhưng ta không muốn để kỹ năng của ta bị thui chột, thế là ta dấn thân vào-''

Tôi nheo mắt, cắt ngang ngay trước khi em ấy kịp nói hết câu: ``Hát rong hả?''

Nàng Succubus  lập tức quay phắt sang tôi, trợn mắt: ``\textbf{Rong cái đầu anh ấy! Hát đường phố!}''

Tôi phì cười trước phản ứng của em ấy, trong khi bác tôi và bối tôi cũng chỉ biết cười trừ. Ông bác tôi lắc đầu, đưa tay quệt mồ hôi: ``Có khác quái gì đâu\dots''

``Khác chứ! Hát rong nghe rẻ tiền lắm!'' em ấy bao biện. Trong mắt tôi, việc đi đây đi đó chỉ để hát thì có khác gì để hát rong đâu.

Bác tôi ho khẽ một tiếng, quay lại vấn đề chính: ``Vậy nếu như thế thì có đủ sống không?''

Aloe-chan khoanh tay, vẻ mặt có chút đăm chiêu: ``Chỉ đủ để sống qua ngày thôi.''

Tôi xoa cằm, khẽ nhún vai: ``Ái chà\dots\ Khoai đấy\dots\ Vậy rốt cuộc là làm nhiều việc nhỉ?''

``Chứ sao nữa! Miễn sao em vẫn được hát thì em sẽ chấp tất!''

Cả ba chúng tôi nhướng mày: ``\dots Và phản ứng của mọi người thế nào?''

Tới đây, cô nàng Quỷ dâm như bừng sáng. Đôi mắt xanh lục long lanh đầy tự hào, em ấy vươn vai một cái rồi tươi cười: ``Tất nhiên là rất tốt chứ! Khả quan là đằng khác!''

Tôi khẽ gật gù. Dù gì thì Aloe-chan vốn có giọng hát rất đặc biệt, chỉ cần có cơ hội, em ấy chắc chắn sẽ tỏa sáng. Tôi gõ nhẹ lên bàn, suy nghĩ một chút rồi hỏi tiếp: ``Vậy tức là em làm nghề gì có động đến âm nhạc à\dots\ Cũng hay đấy chứ. Cơ mà trang thiết bị thiếu thốn thế thì em tập kiểu gì?''

Aloe vung tay đầy tự tin, hất tóc ra sau: ``Anh cứ chớ lo! Aloe ta đây biết tất! Em thậm chí có lần đột nhập vào phòng thu âm của một tòa nhà mà!''

Cả ba chúng tôi đổ mồ hôi hột. Bác tôi và bố tôi ngỡ ngàng nhìn nhau, còn tôi thì trợn tròn mắt. Bố tôi bật cười lắc đầu: ``Thế là liều đấy.''

Bác Hijaki-san thì giơ hai tay lên trời, tỏ ý bó tay bất lực với sự cứng đầu của em ấy: ``Đam mê đến mức này thì tao cũng vái.''

Aloe-chan nhún vai, nở nụ cười tinh quái: ``Cứ cái gì liên quan đến hát và kiếm được tiền là ổn thôi! Ta mặt dày lắm, nhiều lúc còn cố xông vào phòng trà cho bằng được mà!''

Tôi lắc đầu, còn hai ông già bên cạnh thì bật cười. Dường như không khí căng thẳng trước đó đã hoàn toàn biến mất, nhường chỗ cho một bầu không khí nhẹ nhõm hơn.

Hijaki đột nhiên vỗ tay cái bốp, như vừa nhớ ra điều gì đó: ``À đấy, tao có nghe lóng ngóng từ quán café bên cạnh, thảo nào cứ ồn ào thế\dots''

Tôi nhíu mày, khó hiểu trước câu nói vu vơ của bác ấy: ``Ý bác là sao ạ?''

Bác tôi chống tay lên đầu gối, hất cằm về phía ngoài cửa sổ, nơi có thể nhìn thấy một góc bảng hiệu quán café ở ngay bên cạnh nhà chúng tôi: ``Cả tháng nay tao thấy cái quán café bên cạnh nhà mình dạo này bát nháo lắm. Mỗi lần đi ngang qua đều nghe tiếng rộn ràng, náo nhiệt lắm\dots\ Tao cứ tưởng họ có sự kiện gì đặc biệt chứ.''

Chưa kịp để tôi suy đoán, cô nàng Succubus đã nhanh nhảu lên tiếng với vẻ đắc chí: ``Đấy là bởi vì ta đã xin được một chỗ ở đó đấy!''

Tôi hơi bất ngờ, nhưng rồi lập tức bật cười: ``Nếu mà `xin' như em nói thì anh lại thấy họ giống như kiểu `thôi được rồi, tôi cho cô vào làm việc đấy!', vì em làm phiền họ vãi cả chưởng.''

Em ấy lập tức phồng má, mắt trừng lên đầy bất mãn: ``Anh nói thế ý anh là gì hả! Em còn tìm được một chỗ ở trọ gần đó luôn ấy!!''

Tôi nhìn Aloe, rồi chỉ biết xoa trán, thở dài ngán ngẩm: ``Bướng bỉnh đến thế là bó-văn-tay rồi\dots''

Cô nàng khoanh tay trước ngực, gật gù với vẻ tự hào: ``Ta còn nhiệt tình giúp họ ở quán nữa kia mà!''

Tôi nghiêng đầu nhìn em ấy, nheo mắt đầy nghi ngờ: ``Ví dụ như?''

``Pha cà phê này, làm bồi bàn này, rồi\dots''

``Nói chung là làm mọi thứ ở đấy luôn?'' tôi phẩy tay, hiểu được ý của em ấy.

Aloe-chan búng tay một cái đầy tự tin, hất mái tóc hồng ra sau: ``Còn cả hát nữa chứ! Em nói thật, giọng ca của em ấy-''

Không chờ em ấy khoe khoang thêm, tôi liền đưa tay lên chặn lại, chẹp miệng một tiếng: ``Nãy giờ anh thấy em cứ khoe giọng hát của mình với ba người bọn anh, bây giờ em thử hát cho anh nghe xem nào?''

Bầu không khí trong phòng bỗng dưng chùng xuống một chút khi tất cả ánh mắt đều đổ dồn về phía cô nàng Succubus. Cô nàng tóc hồng chớp mắt một cái, rồi môi khẽ cong lên đầy thích thú. Cảm giác như em ấy vừa chờ đúng câu này từ lâu rồi.

Aloe-chan gần như không cần suy nghĩ mà gật đầu cái rụp, sự hào hứng bộc lộ rõ trong từng đường nét trên khuôn mặt: ``Được thôi! Rồi các người sẽ thấy được sức mạnh giọng hát của Mano Aloe ta đây!''

Chúng tôi nhìn nhau, hơi bất ngờ trước sự tự tin thái quá của em ấy. Nhưng cũng chính cái phong thái đó khiến em ấy có một sức hút rất riêng.

Nàng Succubus nhẹ nhàng duỗi đuôi ra phía trước, để lộ một chiếc mic nhỏ nhắn đang cuốn quanh đó. Với một động tác gọn gàng, em ấy tháo ra, lắc nhẹ chiếc đuôi rồi đưa mic lên miệng. Em hắng giọng một hồi, khuôn mặt tỏ vẻ nghiêm túc hơn hẳn.

Và rồi, giai điệu quen thuộc vang lên.

``Ochame Kinou'' - một trong những bài hát mà tôi biết rất rõ trong nội bộ \hll\ - được coi là bài hát đầu tiên của bất kỳ thành viên nào phải hát khi đặt chân tới đây.

Nhưng lần này, nó không phải là giọng của nhóm \hll\ như thường lệ, mà là một phiên bản riêng biệt do chính Mano Aloe-chan thể hiện.

Ngay khi những nốt đầu tiên được cất lên, tôi và bố đều có chút sững sờ.

Một chất giọng trong trẻo, mang đậm sự tươi trẻ. Mỗi câu hát đều chứa đầy nội lực, không phải theo kiểu gồng ép, mà là sự tự nhiên, thu hút người nghe từ những giây đầu tiên. Một giọng hát tinh nghịch, đúng với bản chất quỷ nhỏ của em ấy, nhưng vẫn có nét cuốn hút như một thần tượng chuyên nghiệp.

Giọng hát đó, nếu đem ra so sánh, có thể đặt ngang hàng với Hoshimachi Suisei-chan - người có kỹ thuật đỉnh cao trong Tam Giác Xanh, hoặc thậm chí cả Tokino Sora - thần tượng thuần khiết nhất mà công ty ấy từng có.

Tôi liếc mắt sang bố tôi. Ryujin-san, một audiophile chính hiệu, người đã dành cả đời để sưu tầm và nghiên cứu âm nhạc, cũng đang nhắm mắt, nhịp nhịp ngón tay theo từng giai điệu. Biểu hiện này ở bố tôi không nhiều, chỉ xảy ra khi ông thực sự đánh giá cao một giọng hát.

Aloe-chan kết thúc bài hát với một nốt cao hoàn hảo, rồi thả lỏng người, nhìn chúng tôi đầy kiêu hãnh: ``Thế nào?''

Hai bố con chúng tôi - một audiophile, một Tổng quản lý làm việc ở một công ty đào tạo thần tượng - đều cùng đồng loạt gật gù, đồng thanh đáp: ``Rất tốt đấy.''

Tôi suy nghĩ một chút rồi nhận xét: ``Nếu mà so sánh ra thì em có thể dư sức đấu được với Sui-chan và Sora-chan đấy. Em có một cái giọng đặc biệt khác hẳn với mọi người trong công ty luôn.''

Bố tôi cũng gật đầu tán thành, trầm giọng: ``Chất giọng của cháu có nói lên một sự trẻ trung trong đó, chú phải thừa nhận. Cháu còn hát rõ chữ nữa, lấy hơi đúng kiểu. Chắc cháu phải tập luyện lâu rồi nhỉ?''

Cô nàng lập tức khoanh tay, cười một cách đầy tự mãn: ``Chuyện! Đúng quá luôn! Không có gì phải chê cả!''

Tôi và bố liếc nhau, rồi gần như đồng thời mở miệng:

\begin{dialogue}
    \speak{Kuon} À, cái đó thì\dots
    \speak{Ryujin} \dots cũng không phải là không có điểm yếu.
\end{dialogue}

Nghe chúng tôi nói vậy, em ấy giật thót, đôi tai nhọn hoắt vểnh lên vì bất ngờ: ``\dots Hả?''

Có vẻ như em ấy đang mong đợi một lời khen nữa, chứ không phải một lời góp ý. Nhận ra bầu không khí đang dần thay đổi, tôi quyết định phá vỡ sự im lặng: ``Nhưng mà, nếu xét về tổng thể thì vẫn có một số điểm em cần cải thiện.''

Em ấy hơi khựng lại. Tôi tiếp tục: ``Giọng em hơi bẹt, lại có chút chói khi lên nốt cao. Bố có thấy vậy không?''

Bố tôi gật gù, trầm ngâm một chút rồi nhận xét: ``Hình như chú thấy cháu có vẻ hơi năng động quá hay sao mà giọng cháu thỉnh thoảng lại cục một vó như thế?''

Aloe-chan chớp mắt, như thể không tin vào tai mình. Cô nàng vốn quen với việc tự nâng tầm giọng hát của mình, nhưng khi bị góp ý, phản ứng đầu tiên của em ấy lại là tức tối: ``C-Cái quái\dots\ Sao hai người cứ tìm yếu điểm của ta thế!?''

Bác Hijaki-san khoanh tay, bình thản đáp: ``Tao không nghĩ họ vạch lá tìm sâu đâu, họ đang góp ý đấy. Mày cũng phải tiếp nhận ý kiến để khắc phục chứ.''

Tôi nhún vai, cố giữ giọng điệu trung lập: ``Bọn anh không cố ý làm phật lòng em thật mà. Giọng hát nói chung hay, nhưng vẫn phải cải tạo thêm nữa.''

Bố tôi gật đầu đồng tình: ``Chú cũng phải đồng ý với thằng Cò.''

Nàng Succubus trợn tròn mắt, vẻ mặt hệt như một con mèo bị đổ nước lạnh vào đầu. Giọng điệu của em ấy bỗng chùng xuống, mang theo chút hy vọng mong manh: ``H-Hai người không cảm thấy rung động gì trong lòng hay sao\dots?''

Tôi và bố nhìn nhau, rồi đồng thanh đáp, không chút do dự: ``Chưa đến mức như thế.''

Khoảnh khắc đó, tôi có thể thấy rõ ràng sự sụp đổ trên khuôn mặt Aloe-chan. Nếu như trước đó em ấy còn đầy tự tin và kiêu hãnh, thì bây giờ, gương mặt đó đã méo xệch như một đứa trẻ bị lấy mất kẹo.

Sau khi vừa được nghe những lời nhận xét như ``rót mật vào tai'', để rồi đột ngột bị ``dội gáo nước lạnh'' vào cái tôi của mình, nàng Succubus bật chế độ ăn vạ ngay tại chỗ.

Em ấy thả rơi chiếc mic-rô xuống sàn, rồi nằm lăn ra, đập tay đập chân như một đứa trẻ đang giận dỗi: ``Không chịu đâu, không chịu đâu! Tại sao lại thế chứ!? Rõ ràng là ta hát hay lắm mà!''

Tôi chán nản nhìn Aloe-chan đang lăn lộn trên sàn, cẩn thận lùi lại một bước để tránh bị ảnh hưởng bởi đợt ``bão mè nheo'' này, từ cả chân đạp, tay đấm và đuôi quất: ``Có người nọ người kia em à- Mà cẩn thận lại đạp vào mấy cái hộp bây giờ!''

Bố tôi thở dài, lắc đầu: ``Chú nghĩ là cháu nên học cách chấp nhận chỉ trích đi\dots''

Bác Hijaki-san đứng cạnh cũng chêm vào: ``Công nhận.''

Nhưng tất cả những lời nói đó đều vô dụng, vì em ấy vẫn tiếp tục than vãn. Em ấy ôm đầu, lăn qua lăn lại trên sàn, mè nheo suốt cả mười phút không ngừng nghỉ.

Tôi và bố chỉ có thể bất lực nhìn nhau, trong khi bác Hijaki thì lấy một ngụm trà, lắc đầu ngao ngán.

Tất cả chúng tôi đều cùng chung một suy nghĩ khi thấy cảnh tượng trước mắt: ``Con quỷ nhỏ này đúng là không chịu nổi lời chê thật\dots''

\ct

\begin{center}
    ({\brian Music used: Dark Ages - Loonboon})
\end{center}


Bốn người chúng tôi tiếp tục sắp xếp lại phòng ốc cho Aloe-chan trong khi trò chuyện với nhau. Thậm chí cả vị khách này cũng phải ``động tay động chân'' phụ giúp, vì nếu không, chúng tôi sẽ để mặc em ấy tự dọn dẹp đồ đạc một mình - hoặc ít nhất, tôi dọa em ấy như vậy.

Trong lúc đang lắp đặt máy tính để bàn, tôi quay sang cô nàng hỏi: ``Mà này, em vẫn chưa trả lời câu hỏi tại sao em lại chuyển về đây sống cùng bọn anh đấy.''

Đó cũng chính là câu hỏi đầu tiên mà tôi đã đặt ra cho em ấy, nhưng dường như em ấy đã lòng vòng khi những gì chúng tôi nói chuyện là những khó khăn mà em ấy gặp phải sau sự cố tai hại đó.

Nhưng rồi cuối cùng, điều gì tới cũng phải tới. Aloe-chan đang loay hoay sắp xếp quần áo trong tủ, quay lại đáp: ``Như em vừa nói đấy, YAGOO bảo em dọn về chỗ này. Nhưng tự dưng ông ấy gọi cho em, hỏi em mấy câu rồi đùng một cái, bảo em sống cùng với các người.''

Bác Hijaki-san cau mày: ``Cụ thể là ra làm sao?''

Tôi cũng cảm thấy khó hiểu, liền hỏi thêm: ``Ông ấy đã hỏi gì vậy?''

Cô nàng ngừng tay một chút, lẩm bẩm: ``Chuyện là\dots''

\pov{Mano Aloe}

Chuyện xảy ra khá là dài.

Hôm đó là một ngày bình thường ở quán café Domino, ngay cạnh nhà của cậu Tổng. Ta đã làm việc ở đây được khoảng một tháng và đang tính tiền cho những vị khách cuối cùng trước khi tạm nghỉ trưa.

Ta phải công nhận, bản năng succubus của ta đúng là bá đạo. Nhờ có nó, quán café này hút khách hơn hẳn. Kết hợp với cách nói chuyện tinh nghịch và giọng hát trời cho, ta đã khiến quán làm ăn phát đạt, vượt qua giai đoạn khó khăn do đại dịch. Giờ đây, quán café Domino không chỉ nổi tiếng ở khu phố mà thậm chí còn có tiếng trong cả cái quận Kyuukami (ta nhớ tên của nơi này gọi như vậy).

Kể cả quản lý quán cũng phải ngạc nhiên với kết quả ngoài mong đợi. Chẳng ai ngờ rằng, từ khi một cựu idol của Gen 5 ``xin'' làm việc, doanh thu lại tăng vọt đến thế.

Vừa hoàn tất đơn hàng cuối cùng, ta cười tươi rói, giơ hóa đơn ra trước mặt khách: `` Đơn của các người hết 2850 yên nha\~{}''

Một trong số mấy vị khách nam trẻ tuổi hào hứng đáp: ``Đây, bọn anh bo hẳn 3000 yên luôn, khỏi cần tiền thừa!'' và đưa cho ta hai tờ tiền.

Ái chà, khách sộp đây.

Ta nhướng mày đầy thích thú, chống cằm, giọng điệu kéo dài: ``Hay nha, hay nha\~{} Quý hóa ghê nha\~{} Cho ta xin! Ta tặng cho các người nụ hôn này!''

Rồi ta chu môi, thổi một nụ hôn gió về phía họ. Chỉ vậy thôi mà đám người đó đã bủn rủn tay chân, reo hò ầm ĩ như được ban phước từ trời. Chúng cảm ơn rối rít rồi hí hửng lên xe ra về, trông hệt như một lũ ngốc chính hiệu.

Đợi bọn họ đi khuất, ta khoanh tay, hừ nhẹ một tiếng, vừa lau dọn bàn ghế vừa lầm bầm: ``Thật là, đúng là bọn nhân gian ngu ngốc. Chưa gì mà đã bị ta hạ đo ván dễ dàng rồi\dots''

\begin{center}
    \uline{Tạm ngắt hồi tưởng.}
\end{center}

Kuon-kun khoanh tay, nhướng mày nhìn ta, giọng điệu đầy châm chọc: ``Chuẩn kiểu succubus luôn, thế mà nãy còn chối.'' Anh ấy quay về người bác có khuôn mặt hằm hằm kia: ``Bác Hijaki-san, quán café bên cạnh lúc nào cũng như thế này ạ?''

Hijaki nhăn mặt, bực bội đáp: ``Tao cũng chả biết. Tao chỉ thấy bên kia cứ tưng bừng lắm, hết cả ngày rồi đêm, ung cả đầu!''

Anh Tổng - hoặc ít nhất, anh tự cho như vậy - gật gù: ``Bảo sao mà mấy ngày gần đây lúc cháu về thấy quán cạnh nhà mình đông vui thế.''

Ta chống nạnh, phồng má, tặc lưỡi: ``Các người chưa nghe fan-service bao giờ à? Ta nghĩ các người cũng thường làm thế với khách hàng của mình, nên ta học theo thôi!''

Ryujin chép miệng, lắc đầu: ``Nhưng mà đến mức còn náo loạn như bác Hijaki nói thì chú còn thấy sốc đấy. Chẳng hiểu giới trẻ bọn cháu hút khách kiểu gì nữa\dots''

Kuon-kun nhún vai: ``Đến Gen Z như con cũng bó tay bố ạ.''

Ta hắng giọng, quắc mắt nhìn cả bọn, giọng gắt gỏng: ``*e hèm!* Để ta kể tiếp đây!''

\begin{center}
    \uline{Tiếp tục hồi tưởng.}
\end{center}

Ta vẫn còn đang lau dọn quán cùng với mấy nhân viên khác thì bỗng nhiên điện thoại rung lên.

``Hử? Giờ này ai gọi mình nhỉ?''

Ta với tay lấy điện thoại từ túi tạp dề, liếc nhìn màn hình\dots\ và lập tức sững người.

``\dots Ể?''

Tên người gọi hiện trên màn hình khiến ta đơ mất vài giây.

``YAGOO\dots?''

Cái quái gì đây? Sao ông ta lại gọi cho ta? Rõ ràng là ta đã rời \hll\ nửa năm trước rồi cơ mà?

Ta cứ đứng chôn chân tại chỗ, mắt dán vào màn hình điện thoại, đến độ một nhân viên làm cùng thấy lạ, liền cất giọng hỏi: ``\vie{Aloe-chan? Sao thế?}''

Giật mình, ta vội vã khua tay lên trời, lắp bắp: ``\vie{K-Không có gì đâu! Chỉ là người nhà của ta gọi điện thôi!}''

\dots Sao? Ta cũng có thể nói được thứ ngôn ngữ ở xứ Kawachi với trình độ giao tiếp cơ bản đấy.

Nghe vậy, đồng nghiệp ta không hỏi gì thêm, chỉ nhún vai rồi tiếp tục công việc.

Còn ta, ta nhanh chóng xin phép quản lý ra ngoài để nghe máy.

Vừa bắt máy, giọng nói quen thuộc từ đầu dây bên kia vang lên: ``Lâu rồi không gặp nhỉ, Mano-san?''

Tim ta khẽ nhói lên một nhịp. Ta nghiến răng, siết chặt điện thoại: ``\dots YAGOO\dots Tôi tưởng tôi đã rời công ty rồi mà? Tại sao lại-''

Tanigo ngắt lời ta, giọng vẫn điềm tĩnh như mọi khi: ``Bình tĩnh đi nào. Tôi chỉ định hỏi thăm sức khỏe và công việc của cô thôi, sao phải nóng thế?''

Ta ấp úng, không biết phải trả lời thế nào. Trong lòng dấy lên đủ thứ cảm xúc hỗn loạn, vừa khó chịu, vừa hoài nghi. Ông ta gọi cho tôi vào lúc này để làm gì chứ?

Nhưng trước khi ta kịp nói gì, YAGOO lại bật cười nhẹ rồi tiếp tục: ``Cứ bình tĩnh mà nói đi. Hãy coi như chúng ta chỉ là hai người bạn lâu ngày không gặp thôi.''

Bạn bè? Ta với ông ta á? Ta mím môi, im lặng, không muốn tiếp tục cái cuộc nói chuyện này nữa.

Tuy nhiên, YAGOO không để sự im lặng kéo dài quá lâu. Ông ta phá vỡ nó bằng một câu nói làm tim ta thắt lại: ``Tôi biết, tai nạn hồi đó thực sự rất đáng tiếc. Đặc biệt là với một giọng ca tuyệt vời như cô, thì đó quả là điều không đáng có.''

Cảm giác khó chịu bỗng trào dâng, ta siết chặt điện thoại, giọng run lên vì tức giận: ``\dots Giờ ông đang định trách móc tôi sao? Tôi đã phải chịu đau khổ đến như vậy rồi, ông còn cười cợt tôi được à?''

``Ấy, ý tôi không phải thế,'' YAGOO vội vàng đính chính. ``Điều tôi muốn nói là, ai cũng có thể mắc sai lầm. Đến cả CEO như tôi còn mắc sai lầm nữa kia mà, đâu chỉ riêng mỗi cô?''

Ta nheo mắt, giọng hơi gằn lại: ``Ý ông là sao chứ?''

Đầu dây bên kia lặng đi trong giây lát. Sau đó, YAGOO thở dài một hơi dài, giọng nói trầm hẳn xuống, mang theo chút u ám mà ta chưa từng nghe thấy bao giờ: ``Tôi đã không thể bảo vệ được các idol của mình. Như cô thấy đấy, cô mắc tội một, thì tôi có lẽ mắc tội bốn, năm lần như thế.''

Ta đứng sững, không nói nổi lời nào. Ông ta\dots\ đang tự nhận lỗi sao? Vào thời điểm này ư?

``Tôi biết\dots\ Cô mới chỉ có 15 tuổi thôi. Đương đầu trước sóng gió cuộc đời ngay khi vừa bước chân vào ngành công nghiệp idol này hẳn đã rất khó khăn.''

Giọng ông ta chậm rãi, đầy hối tiếc: ``Lẽ ra tôi đã có thể ở đó bao bọc, che chở cho cô. Nếu như vậy, thì có khi cô đã có thể sát cánh với Omaru-san, Yukihana-san, Shishiro-san, Momosuzu-san, cùng với mọi người ở \hll\ rồi.''

Ta cắn môi. ``\dots Kể cả ông nói thế thì bây giờ cũng có thể thay đổi được gì đâu\dots\ Dù cho ông đã hỗ trợ cảnh sát tống giam gã đó\dots''

YAGOO ngắt lời ngay lập tức: ``Hắn ta không đáng phải bận tâm. Chuyện gì đã qua rồi thì cứ để cho qua thôi. Tôi biết cô đã tự dằn vặt chính mình sau sự cố đó, và có thể cô cũng chẳng còn trông mong gì đến \hll\ cả, nhưng\dots''

Ta nheo mắt, chột dạ: ``Nhưng gì? Dù ông có xin thế nào thì tôi đâu còn có thể quay lại đó được nữa\dots''

``\emph{Sao lại không?}''

Câu nói của ông ta khiến ta chết lặng. ``Ể?''

YAGOO tiếp tục: ``Tôi cũng biết là cô đã thử làm nhiều nghề khác nhau, và giờ đang có công việc ổn định ở quán café ở Kawachi rồi nhỉ?''

Cơn sốc nối tiếp cơn sốc.

Ta bật thốt lên: ``Ủa!? Sao ông biết!?''

YAGOO cười nhẹ, đáp đơn giản: ``Tôi có cách riêng của mình thôi. Nhưng mà hỏi nhanh đáp gọn này, thu nhập cũng ổn chứ?''

Ta nhún vai, mặc dù biết rằng ông ta chẳng thể nào thấy được: ``Đủ để trả tiền trọ hàng tháng. Mà cái này có liên quan gì đến ông? Tôi có làm việc với ông nữa đâu?''

YAGOO không vội trả lời ngay. Đến khi ông ta cất giọng, giọng điệu ấy không còn là của một CEO, mà giống như một người anh lớn đang khẽ nhắc nhở đứa em gái bé bỏng của mình: ``Liên quan chứ. Dù gì cô cũng từng là một thành viên, và đối với mọi người, cô vẫn sẽ luôn mãi là thành viên của \hll.''

Tim ta như bị bóp nghẹt.

Miệng ta hé ra, định nói gì đó, nhưng rồi lại chẳng thể thốt lên lời.

Im lặng. Ta vẫn chưa thể tiêu hóa hết câu nói trước đó của YAGOO.

Không đợi ta kịp phản ứng, ông ta tiếp tục: ``Cô cũng đang tập luyện chăm chỉ để hướng tới trở thành một diva nơi xứ Quỷ, đúng chứ? Tôi xem các video utaite trên mạng rồi, và nói thật, cô cũng có cải thiện chút ít đó.''

Lần này, ta không kìm được mà bật cười nhạt. ``\dots Tôi cúp máy đây.''

``Ấy từ từ, đừng cúp!'' YAGOO vội lên tiếng, giọng có chút khẩn trương, nhưng ta vẫn giữ nguyên tư thế bấm nút kết thúc cuộc gọi.

``Hay là thế này, tôi sẽ cho cô một ưu đãi, nếu không được thì ta coi như từ mặt nhau, được chứ?''

Nghe đến đây, ta hơi nhíu mày. ``\dots Ông nói đi cho tôi nghe thử xem nào.''

Ông ta hắng giọng một chút, rồi chậm rãi cất lời: ``Được rồi, như thế này\dots''

Ta bỗng cảm thấy có gì đó kỳ lạ trong giọng điệu của ông ta, như thể ông ta đang thử thách ta.

``Cô để ý đến cửa hàng viễn thông ngay bên cạnh quán café cô đang làm chứ?''

Theo phản xạ, ta nhìn sang bên trái. Đúng là có một cửa hàng viễn thông đang làm việc bình thường trong những ngày giáp Tết, ánh đèn từ bảng hiệu phát ra sáng rực một góc phố.

``Có. Sao nữa?''

``Đó là nhà của một trong những nhân viên kỳ cựu của công ty \hll\ chúng ta. Cô biết cậu ta chứ?''

Câu hỏi làm ta hơi ngạc nhiên. ``Không. Sao?''

``Khi nhóm \textit{holofive} họp online từ tháng Bảy năm ngoái cùng với chính người đó thì cô lại không có mặt, nên chắc cô không biết đâu.''

Ta nhướng mày, giọng có chút khó chịu: ``Hôm đó tôi ốm, không thể dậy được đấy thôi. Cái đó thì liên quan gì?''

``Họ không nói gì với cô sau đó à?''

``Tôi chỉ nhớ họ nói rằng đó là một cậu tầm 24, 25 tuổi thôi. Nghe đồn là cậu ta được mọi người yêu quý lắm.''

Ta không nhớ rõ, nhưng hình như Nene có kể qua về người đó, một lần duy nhất, trong tin nhắn thoại kéo dài chưa đến ba mươi giây.

``Đúng rồi, đó chính là nhà của cậu ta đấy,'' ông ấy xác nhận. ``Sắp tới chúng tôi sẽ cải tạo nó thành ký túc xá dành cho những cựu thành viên \hll\ như cô có thể sinh sống ở đó.''

Ta trố mắt. ``\dots Thế thì liên quan gì đến tôi?''

\begin{center}
    \uline{Tạm ngắt hồi tưởng.}
\end{center}

``Đến đây thì anh biết cái vụ ``ký túc xá'' ở đâu ra rồi đấy\dots'' Kuon-kun gật gù, cắt ngang lời nói của tôi.

``Ủa? Anh đã biết rồi sao?'' tôi nghiêng đầu. ``Anh luận được rồi sao?''

``Suy luận logic. Về cơ bản lại là cái trò gì đó điên rồ nữa của CEO chúng ta mà anh còn không biết cơ,'' Kuon-kun đáp lại.

``Ông ta luôn làm cho mọi người ở công ty nó lẫn cả gia đình chú bất ngờ đấy,'' Ryujin tiếp lời ông ấy, ngán ngẩm nói. ``Tết năm ngoái nguyên cả công ty sang nhà chú để quay phim đấy.''

``Ông ta luôn có nhiều chiêu trò quái thai đến thế ư?'' ta nhíu mày.

``Cũng không đến mức như thế. Nhưng nếu làm những thứ kiểu này thì ắt phải có âm mưu gì đó, chứ anh biết là Tanigo-san sẽ không làm một điều gì đó ngớ ngẩn đến mất não đâu,'' anh Tổng - tự nhận - đáp lại. Giờ ta mới biết là nhân viên thậm chí còn nghĩ xấu cả sếp đấy.

``Vậy là mày đồng ý chuyển về đây sống à?'' Hijaki cất giọng hỏi.

``Ban đầu ta cũng do dự lắm, kiểu tại sao ta phải sống cùng với các người, trong khi ta chẳng còn dính líu gì đến \hll\ nữa, nhưng rồi\dots''

Ta chợt ngừng lại, cố gắng tìm ra một ý để tiếp tục trình bày.

``\dots Gì nữa?''

Ta thở dài, chẹp miệng: ``\dots Để ta kể tiếp đây.''

\begin{center}
    \uline{Tiếp tục hồi tưởng.}
\end{center}

Ta đã hỏi vị CEO kia với giọng ngờ vực: ``Tại sao tôi phải ở cùng với họ?''

Nhưng thay vì trả lời thẳng, ông ta lại hỏi ngược lại: ``Cô có muốn được gặp lại bạn bè của mình không?''

Tim ta khẽ hẫng một nhịp.

Không cần ông ta nói thêm, ta cũng hiểu câu hỏi ấy có ý nghĩa gì.

Ta im lặng. Không phải vì ta không muốn trả lời, mà vì ta không dám trả lời.

Như thể nhìn thấu mọi suy nghĩ trong đầu ta, YAGOO tiếp tục: ``Trong một buổi collab duy nhất của cô với Gen 5 ấy, mọi người trông ăn ý với nhau lắm đấy. Chứng tỏ rằng cô cũng đã làm bạn với họ từ lâu rồi nhỉ?''

Ta siết chặt điện thoại trong tay. Hồi ức ngày đó chợt ùa về.

Buổi phát sóng chung duy nhất của năm người bọn ta ngày hôm đó. Một buổi phát sóng với nhiều tiếng cười, những trò đùa, và cả niềm vui nữa.

Đó cũng là lần đầu tiên, và là lần duy nhất, bọn ta có được những giờ phút hạnh phúc tới vậy.

``\dots Bọn tôi đã làm bạn từ sáu năm nay rồi. Những lúc vui buồn có nhau, mỗi khi chơi điện tử lại rủ nhau sang nhà Botan, rồi còn\dots''

Giọng của ta nhỏ dần.

Ông ta không nói gì một lúc lâu, như thể cố tình để ta tự ngẫm nghĩ.

Rồi, bằng một giọng trầm hơn hẳn lúc trước, ông ta khẽ nói: ``Tôi hiểu mà. Nhưng\dots\ tại sao cô lại không liên lạc với họ? Đến Momosuzu-chan đã cố gắng gọi cho cô nhiều lắm đấy.''

\dots Nene? Cậu ta\dots\ đã gọi cho tôi sao?

``\dots Làm sao mà tôi có thể gặp lại bạn mình khi tôi bị loại khỏi vòng gửi xe chứ?'' tôi nói, cố kìm nén cảm xúc mình lại. ``Bọn họ đã là idol rồi, trong khi tôi vẫn phải chật vật kiếm sống\dots\ Ai cho tôi tự trọng chứ? Nếu có cho thì tôi cũng không làm được!''

``\emph{Vì họ cũng nhớ cô lắm đấy.}''

Ta khựng lại. Những người đồng đội đó\dots\ thực sự lại thèm nhớ tới một kẻ bất tài như ta sao?

``Đừng cứng đầu nữa. Tôi biết cô cũng muốn gặp lại bọn họ, từ sâu trong đáy lòng kia mà.''

Ta lắp bắp, cố gắng phản biện lại: ``T-Tôi không\dots''

Nhưng rồi, ông ta lại tiếp tục: ``Tôi là CEO của \hll\ mà, tôi cũng học chút ít từ cậu nhân viên chăm chỉ kia, nên tôi thấu hiểu rất rõ mọi người, kể cả cô nữa.''

Làm sao mà ta có thể gặp lại họ được?
Làm sao ta có thể nhìn thẳng vào mắt họ, khi ta là người đã rời bỏ tất cả, khi chính ta đã tự mình bước ra khỏi vòng tròn ấy?

Không phải vì ta không muốn liên lạc.
Không phải vì ta ghét họ.

Chỉ là\dots\ ta sợ.

Ta sợ ánh mắt đầy lo lắng của họ.
Ta sợ những câu hỏi mà ta không thể trả lời.
Ta sợ cảm giác rằng mình đã đánh mất vị trí bên cạnh họ, và giờ đây ta chỉ còn là một kẻ ngoài cuộc.

Ta muốn chạy trốn khỏi thực tại.
Vậy mà\dots

``Kể từ khi cô rời công ty không kèn không trống, bốn người họ đã cố gắng gọi cho cô, rồi nhắn tin nữa, nhưng rồi vẫn bặt vô âm tín.''

``\dots Họ lo cho một đứa mắc phải sai lầm như tôi sao?'' tôi nói, giọng run run.

``Đừng tự trách bản thân nữa. Họ là bạn của cô từ nhiều năm rồi, hiển nhiên họ phải lo chứ! Họ có quan tâm mấy chuyện kiểu cô đang ở vị trí nào đâu?''

Trong lòng ta giờ đây là một mớ hỗn độn.

Ta\dots\ có quyền được nhận sự lo lắng ấy sao?

``Momosuzu-chan thậm chí đau buồn đến mức đổ bệnh cả tháng trời đấy, may mà hồi phục được. Tầm hai, ba tháng nữa là có thể đến công ty ra mắt rồi. Cơ mà, những người khác cũng nhung nhớ tới cô lắm chứ.''

Mắt ta đỏ hoe, những giọt nước mắt bắt đầu rơi lã chã.

Tại sao chứ?

Dù ta đã không còn ở \textit{holofive} nữa, nhưng\dots

\dots họ thực sự vẫn quan tâm, vẫn đau buồn vì một người như ta ư?

``\dots Thật à\dots?'' ta vừa nhíu mày ngạc nhiên, vừa ngỡ ngàng với những gì ta nghe được.

YAGOO tiếp tục, giọng chắc nịch: ``Tôi không nói phét đâu. Kể từ tai nạn đó, tôi đã tự kiểm điểm bản thân, đã cố hết sức để bảo vệ tất cả mọi người trong công ty rồi.''

``Vậy\dots\ họ vẫn ổn chứ?'' ta cất giọng. Dẫu biết rằng ta đã rời khỏi vòng tròn kia, ta vẫn không thôi lo lắng với họ.

Ông ta đáp lại: ``Dĩ nhiên. Kể từ sau vụ việc, bốn người bọn họ càng khăng khít hơn hẳn, thậm chí còn bảo vệ lẫn nhau để không mắc lại sai lầm đấy.''

Tôi thở phào nhẹ nhõm: ``Tốt quá rồi còn gì. T-Tôi hy vọng là mọi người ở đó vẫn ổn\dots''

``Còn hơn cả ổn nữa chứ! Nhưng dù bề ngoài họ vui vẻ thế nào, trong thâm tâm, họ biết rằng nhóm \textit{holofive} sẽ không còn rôm rả như ngày trước nữa,'' ông ta nói, giọng tỏ vẻ tiếc nuối. ``Shishiro-chan với Yukihana-chan còn nói với tôi và Hikawa-san rằng họ muốn gặp lại cô nếu có cơ hội mà.''

Ta chợt im lặng một lúc lâu, khựng lại trước những câu từ mà ông ấy nói.

Họ thực sự\dots\ muốn gặp ta ư?

Nhưng, liệu điều đó có còn khả thi không? Ta đã không liên lạc với họ kể từ khi ta rời nhóm, nên không đời họ biết ta đang ở đâu, đúng chứ?

Nhưng, ta cảm tưởng rằng những lời mà YAGOO nói không phải là những lời nói dối, cũng không phải là những lời nói suông chỉ để thỏa mãn ta.

Không đời nào ông ta lại lấy cớ bạn bè để trêu ngươi với ta, đúng chứ?

Lần này, ta không thể phủ nhận nữa. Ta biết rõ câu trả lời trong lòng mình.

``Bây giờ, tôi hỏi lại. Cô đã biết tình hình hiện tại của nhóm các cô rồi, vậy cô có muốn chấp nhận lời đề nghị này không?''

Ta nhắm mắt lại. Hít một hơi thật sâu.

Nếu cho sau này ta có phải làm quần quật, hay thậm chí là ở một xó nào đó đầy xa lạ\dots

\dots nhưng ta vẫn có thể gặp lại bạn bè của ta, dù chỉ là một lần\dots

``\dots Được rồi, tôi chấp nhận.''

Ngay lúc tôi nói ba từ đó - dù có vẻ hơi lưỡng lự - YAGOO thở dài nhẹ nhõm, giọng vui tươi hẳn khi tôi đồng ý với đề nghị này.

Bầu không khí giờ đã bớt căng thẳng hơn. Ông ta nói tiếp: ``Căn nhà của cậu con trai đó, tôi nói trước, là nơi mà nhiều thành viên của công ty tôi hay lui tới để thăm hỏi đấy, theo lời cậu ta báo cáo như vậy.''

Ta nghiêng đầu thắc mắc, mặt hơi co lại vì vừa tò mò, vừa lo sợ: ``Vậy anh ta là ai?''

Và rồi, ông ấy trả lời ráo hoảnh: ``Hãy tìm người có tên Hikawa Kuon. Cậu ta sẽ giúp cô thích nghi với môi trường sống này đấy. Thế nha! ''

Nói xong, chẳng nói chẳng rằng mà cúp máy ngay lập tức, khiến ta không kịp phản ứng gì.

Ta nhìn vào chiếc điện thoại trong tay một lúc lâu, trong đầu chất đầy những dấu hỏi lớn:

``Liệu ông ta có lừa mình không?''

``Cậu Kuon này là thằng nào?''

``Mình sẽ sống ở đâu?''

Ngay lúc ấy, một tin nhắn từ YAGOO gửi đến, như thể ông ta đã đoán trước được những thắc mắc của ta:

\begin{tcolorbox}[
        colframe=vol,    % Màu viền
        colback=white,   % Màu nền
        boxrule=0.5pt,   % Độ dày viền
        arc=0mm,         % Góc vuông (không bo tròn)
        left=5pt,        % Khoảng cách giữa nội dung và viền trái
        right=5pt,       % Khoảng cách giữa nội dung và viền phải
        top=5pt,         % Khoảng cách giữa nội dung và viền trên
        bottom=5pt       % Khoảng cách giữa nội dung và viền dưới
    ]

    Yên tâm đi, nói lời rồi, tôi giữ lời hứa tốt lắm. Cứ chuyển dần đồ đạc sang tầng hai, phòng bên trái đi, rồi tôi sẽ nói chuyện với chủ nhà bên đó. À mà nói trước, cô sẽ chưa thể gặp các thành viên \hll\ bây giờ được đâu, nhưng nghe tôi nói: cứ giữ bí mật đến khi nào thời điểm thích hợp.

    À, và tiền trọ thuê nhà cứ để công ty lo. Cô không cần phải lăn tăn nữa hen!

    \hfill\textbf{-YAGOO.}
\end{tcolorbox}

Ta đọc tin nhắn mà càng thấy khó hiểu hơn.

Tại sao ông ta lại làm đến mức này cho ta chứ?

Chỉ là một cựu thành viên đã rời đi từ lâu\dots\ vậy mà ông ta vẫn quan tâm sao?

Cái đuôi của ta khẽ quẫy nhẹ, theo bản năng gãi đầu trong vô thức.

Trước khi kịp nghĩ thêm, quản lý quán gọi ta quay lại làm việc.

``Thật tình\dots\ Tại sao ông ta lại tốt với mình như vậy chứ?''

Dù có suy nghĩ bao nhiêu đi nữa, ta cũng không thể hiểu được dụng ý của vị CEO \hll\ kia. Nhưng có lẽ, thay vì đặt ra những câu hỏi không có lời giải, tốt nhất ta nên làm theo những gì ông ta đã sắp xếp.

Và thế là, ta dọn đồ từ khu nhà trọ cũ cách chỗ làm không xa, chuyển đến nơi ở mới - cũng chính là nhà của gia đình Hikawa.

\begin{center}
    \uline{Kết thúc hồi tưởng.}
\end{center}

\pov{Hikawa Kuon}

``Tuyệt vời lắm YAGOO-san. Nước đi kiểu này thì đến bố con thằng này còn không đỡ nổi\dots'' Tôi lắc đầu ngao ngán sau khi nghe câu chuyện từ Aloe-chan.

Quả thực, Tanigo-san có những quyết định mà tôi chẳng bao giờ có thể đoán trước được. Từ việc mời một cựu thành viên \hll\ sang nhà tôi ở, đến việc biến nhà của tôi trở thành một ký túc xá.

``Vậy là mày nghe lời ông ta, rồi dọn đồ đạc chuyển lên đây hả?'' bác Hijaki-san nói, dường như cũng không thể tin nổi.

``Ông ta cho ta đãi ngộ tốt như thế, làm sao mà ta có thể cưỡng lại được! Ta đã không muốn về nhà vì sợ bị mất mặt rồi, nên chỉ còn cách bám víu ở đây đến khi thực hiện ước mơ của ta thôi!'' em ấy đáp lại, cố gắng bao biện lại.

``\dots Chậc. Làm tới bến luôn đấy,'' tôi bình luận.

``Bọn chú tưởng là cháu ở đây ít ngày, chứ đâu có nghĩ rằng sẽ ở ghép đến như vậy đâu! Lại còn có cả vụ ký túc xá nữa!'' bố tôi lên tiếng, vừa bày tỏ ngạc nhiên, vừa bày tỏ sự bất bình.

``Ông sếp của mày cũng quái dị không kém đâu, Kuon ạ.''

Tôi thở dài đáp lại: ``\dots Cháu cũng chẳng biết nên coi đó là lời khen hay chê nữa ạ.''

Aloe-chan lên tiếng, dường như cũng vừa trải qua cảm giác ngỡ ngàng và khó tin như vậy: `` Ta cũng thấy ngạc nhiên như các người vậy đó! Ông ta bảo rằng ta sẽ phải sống ở đây một thời gian dài đấy.''

Tất cả ba người đàn ông chúng tôi cùng nhau đồng thanh: ``Như thế là nhà mình có hàng xóm mới còn gì.''

Chúng tôi tiếp tục bàn chuyện rôm rả với nhau trong khi sắp xếp đồ đạc cho vị khách mới của nhà mình.

Chẳng ai trong gia đình tôi nghĩ rằng sẽ có ngày nhà Hikawa biến thành ký túc xá \hll\ theo ``đơn đặt hàng'' của chính CEO công ty cả.

Thật lòng mà nói, tôi không biết dụng ý của Tanigo-san là gì.

Có thể ông ấy chỉ đơn thuần muốn cho Aloe-chan một cơ hội làm lại, nhưng cảm giác như có gì đó sâu xa hơn.

Tôi dám chắc rằng nếu A-san biết chuyện này, chị ấy cũng sẽ bất lực mà thở dài như tôi thôi.

\ct

\begin{center}
    ({\brian Music used: Dark Ages - Demonstration Mini-Game})
\end{center}

Trong lúc chúng tôi hoàn thiện những bước cuối cùng của căn phòng mới cho Aloe-chan, cả nhà bận rộn chia nhau từng việc.

Người xếp cốc, người gấp quần áo, người lắp đặt nội thất.

Cuối cùng, tất cả lại cùng nhau lau dọn lần nữa, đảm bảo từng góc nhỏ đều sạch sẽ không tì vết.

Tầng hai ban đầu vốn bị chia làm hai phần, một nửa - nơi mà chúng tôi đang đứng - thuộc về căn phòng của Aloe-chan, nửa còn lại là không gian chung của gia đình Hikawa chúng tôi. Để tạo thêm sự rộng rãi, tôi và bác Hijaki-san cùng Aloe-chan đã phải phá bỏ bức tường ngăn cách giữa hai bên, ghép hai phần không gian lại thành một.

Trong khi đó, bố tôi cẩn thận sắp xếp từng dụng cụ âm nhạc của nàng Succubus - một công việc không hề nhẹ nhàng khi có đủ thứ từ giá để micro đến mấy cái dây kết nối lằng nhằng. Nhưng dù gì, đây cũng là chuyên môn của ông ấy.

Ông vừa mở một trong những hộp bìa lớn vừa lắc đầu, vẻ mặt không giấu được sự bất ngờ: ``Chú cũng ngạc nhiên khi thấy có cả bộ loa với mixer trong này.''

Bố tôi nói, nhấc lên một trong những thiết bị khá nặng. ``Nhìn chung, cháu cũng có vẻ đầu tư nhỉ?''

Tôi liếc nhìn Aloe, cười xòa: ``Nhìn thậm chí còn xịn hơn nhà mình đấy chứ.''

Cô nàng chỉ nhún vai: ``À, cái đống đó đa phần không phải của ta đâu.''

Tôi liền ngờ ngợ ra ngay lập tức. ``Chắc là YAGOO tặng chứ gì?''

Em ấy gật đầu không chần chừ.

``Từ trên trời rơi xuống luôn?'' tôi hỏi tiếp, và lần này, nàng Succubus lại tiếp tục gật.

Tại sao tôi lại không thấy ngạc nhiên nhỉ? À thực ra là có, nhưng mà tại sao điều đó lại có thể xảy tới cả một người là cựu thành viên của công ty, đã thế lại còn ở ngoài phạm vi thành phố Điện tử này?

Có giời mới biết được.

Bố tôi bày tỏ sự tò mò: ``Thế đồ của cháu mang lên gồm những gì?''

Aloe-chan nghiêng đầu nghĩ ngợi, rồi dùng đuôi vẫy vẫy về phía tôi.

``Chỉ có cái micro cuốn trên đuôi này, cái tai nghe này\dots'' - em ấy chỉ vào chiếc tai nghe trùm đầu đang đeo trên cổ - ``và\dots\ còn gì nữa nhỉ?''

Tôi nhướng mày hỏi tiếp: ``Mang từ quê nhà lên à?''

Aloe-chan liếc tôi một cái, vẻ mặt khó chịu: ``Anh hỏi đểu thế? Không từ nhà thì từ đâu?''

Xong rồi, em ấy nhìn tôi, rướn vai: ``Mà anh trông có vẻ thích thú lắm nhỉ?''

Tôi nhún vai: ``Đồ điện tử thì anh thích mà. Hỏi tí thôi, sao căng?''

Tôi vô thức nhìn chằm chằm vào chiếc tai nghe chụp đầu mà Aloe-chan đang đeo. Hàng có vẻ là hàng xịn đấy, chứ chẳng đùa.

``Bộ anh thấy cái thứ đang đeo trên cổ này còn đáng giá hơn mặt của em sao?'' Câu hỏi đầy vẻ tự ái khiến tôi giật mình. Hóa ra là em ấy đang hiểu nhầm.

``Ờm\dots\ Không hẳn.'' Tôi quay mặt đi chỗ khác, hơi đỏ mặt một chút.

Nhưng thế là chưa xong.

``Đừng có bơ em chứ!'' Em ấy bực mình nhích lại gần, tay chống hông. ``Khuôn mặt hớp hồn của em đây mà anh không có ý kiến gì sao?''

Tôi nhìn em ấy một giây, rồi đáp tỉnh bơ: ``Nói bỗ bã ra thì mặt em nhìn giống con nai lắm. Ngây thơ chưa biết gì.''

Mặt của nàng Succubus đỏ bừng. ``C-Còn ngực thì sao!?''

``Ờm\dots\ Bình thường,'' tôi đáp gọn, dù phải thừa nhận tôi có liếc vào ``cặp bưởi'' của em ấy một ly.

``Vậy còn-''

Tôi cắt ngang trước khi em ấy có thể tiếp tục: ``Tất cả đều ổn cả. Anh thấy cũng thường thôi.''

Mặt Aloe từ đỏ bừng chuyển thành tái mét vì tức giận: ``Ý anh là sao chứ!? Anh thấy cái tai nghe của em còn quý trọng hơn cả ngoại hình của em sao!?''

Bố tôi đành phải chen vào trước khi cuộc cãi vã của hai đứa tôi leo thang: ``Kuon, đừng trêu em ấy nữa. Aloe-chan, thông cảm cho anh ấy đi. Con nhà chú tồ lắm.''

Bác Hijaki-san thì cũng nhịn được, hùa theo: ``Gần ba mươi tuổi đầu rồi mà vẫn không có cảm tình gì với con gái à? Thế này thì hỏng!''

Tôi thở dài, lắc đầu chịu trận: ``Kệ cháu đi ạ.''

Tôi lập tức đổi chủ đề, dường như chưa muốn động đến vấn đề nhạy cảm kia vội. ``Mà nhắc đến đồ em mang sang, ngoài cái micro với cái tai nghe thì em còn mang theo gì nữa?''

Mắt của cô nàng Succubus sáng lên như vừa sực nhớ ra điều gì đó khá quan trọng: ``À đúng rồi! Cả cái loa mini!''

Bác tôi nghiêng đầu, chưa hình dung ra ngay.``Cụ thể là như thế nào?''

Cô nàng Succubus phấn khích giải thích, tay vẽ vẽ trong không khí như đang miêu tả một thứ rất đặc biệt: ``Nó là cái loa mà ta có thể kéo đi khắp nơi! Ngày trước ta cũng dùng nó để đi hát trên đường phố đó!''

Tôi nhìn em ấy, rồi chỉ đáp lại đúng một từ với biểu cảm bất lực: ``\dots OK?''

Aloe-chan nghiêng đầu, khó hiểu trước phản ứng lạnh nhạt của tôi. ``Anh có vẻ không hứng thú lắm nhỉ?''

Tôi thở dài, chống tay lên hông, thở dài đáp: ``Cái loa đó bọn anh gọi là loa kẹo kéo. Và\dots\ ừ. Nó phiền nhiễu lắm.''

``Ể? Sao thế?''

Em ấy chớp mắt, có vẻ vẫn chưa nhận ra vấn đề.

Tôi khoanh tay, nhún vai như thể sắp kể một chuyện nghe có vẻ sẽ động trời, nhưng lại hoàn toàn là sự thật: ``Nói ra không tin, nhưng\dots\ tại nơi anh, nó chính là nguyên nhân xảy ra các vụ án mạng đấy. Đặc biệt là ban đêm.''

Aloe giật mình. ``Thế ư? Nhưng bình thường ta hay luyện giọng vào đêm mà?''

Bố tôi nãy giờ vẫn im lặng, cuối cùng cũng lên tiếng với giọng điệu đầy ẩn ý: ``Và đó chính là lý do chết người đấy.''

Aloe-chan thoáng sững người, nhưng lập tức chống chế: ``T-Ta biết mà! Thế nên dạo gần đây ta chuyển sang luyện lúc sáng rồi! Ta biết nó gây ảnh hưởng trật tự thế nào mà!''

Tôi nhìn em ấy, lần đầu tiên trong ngày thấy được một chút tự giác ở cô nàng kusogaki này.

``\dots Ít nhất với một đứa con gái 15 tuổi như em mà đã biết điều cơ bản như thế là ổn rồi.''

Bác Hijaki đứng khoanh tay, gật gù suy nghĩ: ``Tao nghĩ là phòng này nên lắp thêm hệ thống chống tiếng ồn nữa cho khỏi làm phiền cả nhà\dots''

Tôi gật đầu đồng tình: ``Ý hay đấy bác ạ.''

Nhắc đến nhạc, tôi chợt nhớ ra một chuyện quan trọng mà có thể giúp ích cho Aloe-chan rất nhiều.

``Nhắc đến nhạc thì\dots\ anh nghĩ có một chỗ em có thể làm để nâng cao khả năng của em đấy.''

Em ấy nheo mắt, tỏ vẻ nghi ngờ. ``Gì vậy?''

Tôi nhếch môi cười nhẹ, không trả lời ngay mà chỉ ném lại một gợi ý mở: ``Cái này thì nói chuyện với mẹ anh nhá. Bà ấy biết rõ hơn đấy. Hoặc là bố anh cũng được.''

Nàng Succubus nhìn tôi đầy dè chừng, nhưng cuối cùng cũng nhún vai, tặc lưỡi: ``\dots Em tạm tin anh vậy. Dù gì ta cũng định vãn việc ở quán café rồi.''

\ct

Sau một hồi sắp xếp lại đồ đạc, cuối cùng chúng tôi cũng đã xong. Căn phòng mới của Aloe-chan giờ đây trông rộng rãi và tiện nghi hơn bao giờ hết. Vì đã phá bỏ bức tường ngăn cách giữa hai phòng nhỏ trước đó, không gian này giờ rộng gấp đôi so với các phòng trong nhà.

Một bên phòng được bố trí như một studio thu nhỏ với hệ thống âm thanh chuyên nghiệp: một bộ loa cỡ lớn đặt ngay góc tường, mixer và các thiết bị đi kèm cùng bộ máy tính cây được bố tôi sắp xếp gọn gàng bên bàn làm việc, micro và tai nghe treo ngay ngắn trên giá. Ngay sát đó, một chiếc ghế xoay êm ái cùng với một tấm thảm trải sàn giúp em có thể thoải mái sáng tác hoặc luyện giọng.

Bên còn lại là không gian sinh hoạt cá nhân, bao gồm một chiếc giường đơn kiểu Nhật đặt sát tường, một tủ quần áo nhỏ và một kệ sách lấp đầy bởi những cuốn truyện tranh và sách hướng dẫn về âm nhạc. Một bộ đèn LED viền quanh góc phòng, tạo nên một không gian vừa ấm áp vừa hiện đại, rất phù hợp với phong cách của một Diva Succubus xứ Quỷ giới.

Ba bác cháu tôi cũng đưa ra một vài quy định trong nhà, trong đó có\dots

``Bình thường là em sẽ phải trả tiền thuê nhà ở đây, cơ mà\dots'' tôi mở lời, liếc nhìn Aloe-han xem em phản ứng như thế nào.

Bác Hijaki-san cười khẽ, khoanh tay lại: ``\dots Chắc sếp của mày bao phần này rồi nhỉ?''

Nàng Succubus vẫy tay, vẻ như chuyện đó chẳng có gì đáng lo: ``Ta nghĩ là vậy. Mọi người không phải lo đâu.''

Tôi cũng đoán được trước câu trả lời này. Với một người lắm tiền nhiều của và cũng rất nhiều chiêu như Tanigo-san, việc lo chỗ ở cho em ấy có lẽ chỉ là chuyện nhỏ.

Rồi đến\dots

Bố tôi tiếp lời với giọng điềm đạm hơn: ``Ăn uống thì có thể làm bếp ở đây, hoặc lên nhà ăn cùng các chú.''

Em ấy bĩu môi, gãi gãi má một cách lưỡng lự: ``T-Ta thích ăn một mình hơn. Ta vẫn chưa quen với môi trường này đâu\dots''

Bố tôi gật đầu, không ép. ``Tùy cháu. Cháu không phải ngại đâu.''

Tôi nhún vai, bổ sung thêm: ``Thích thì cứ lên nhà chơi thôi. Nhớ báo sớm một tiếng là được.''

``\dots Dạ,'' em ấy trả lời, giọng vẫn hơi ngập ngừng.

Và cuối cùng\dots

Bác Hijaki cất giọng chắc nịch: ``Công việc dưới này mày tự quản nha. Ryujin với Kuon cũng phải có ý thức trông coi em ấy đấy.''

Aloe-chan lập tức chớp mắt, nghiêng đầu hỏi lại: ``Nghĩa là ta làm gì cũng được à?''

Tôi bật cười, khoanh tay lại. ``Đương nhiên. Cứ tự nhiên như ở nhà thôi. Miễn là đừng ảnh hưởng tới các phòng xung quanh là được.''

Em ấy đập tay lên ngực, đầy tự tin: ``Cái này ta làm tốt! Yên tâm!''

Tôi thầm nghĩ: ``\textit{Cái gì mà 11làm tốt'' chứ\dots\ Cứ chờ xem đã\dots}''

\il

Dù sao thì, cô nàng diva succubus Mano Aloe cũng chính thức trở thành một thành viên tạm thời (?) của gia đình Hikawa chúng tôi.

Không chỉ sống ở đây, em ấy còn có thêm một công việc mới - đó là hát rong tại các hội chợ cùng với mẹ tôi mỗi khi có dịp. Dĩ nhiên, đó không phải là nghề chính, thù lao cũng chẳng đáng là bao và khá thất thường. Nhưng với một người yêu ca hát như Aloe-chan, có lẽ em ấy sẽ không quá bận tâm đến chuyện đó. Ít nhất, đây cũng là một cơ hội để em có thể mài giũa giọng hát của mình một cách thực tế nhất.

Mà nghĩ đi nghĩ lại\dots\ nếu đã thế thì có lẽ tôi nên đầu tư thêm một dàn karaoke đặt dưới phòng em ấy. Aloe-chan thì chắc chắn sẽ thích rồi, nhưng tôi cũng có lý do riêng của mình - với tư cách là đồng chủ nhà, anh đỡ đầu và tổng quản lý của \hll, có lẽ tôi cũng nên tập hát nghiêm túc hơn một chút\dots

Hy vọng em ấy có thể quen với môi trường mới này.

Chào mừng em đến với gia đình Hikawa, Mano Aloe-chan!

\il

Tết Tân Sửu năm nay, gia đình Hikawa có thêm một thành viên mới. Đúng vậy.

Tôi nhớ lại năm ngoái, là Tết Canh Tý, chúng tôi đã tiếp đón hơn hai mươi người tới chung vui. Không khí náo nhiệt, tiếng cười rộn ràng khắp nhà. Nhưng năm nay thì khác, lần này không chỉ đơn thuần là tiếp đón khách trong một buổi tiệc.

Năm nay, chúng tôi chào đón một người sẽ sống cùng gia đình Hikawa một cách lâu dài.

\dots Liệu đây có phải là khởi đầu cho một điều gì đó đặc biệt hơn không nhỉ?

\il

Tôi dựa lưng vào ghế, mắt nhìn chằm chằm lên trần nhà. ``Ký túc xá của nhà Hikawa'', cái danh xưng nghe như đùa ấy đã thành hiện thực chỉ trong chớp mắt.

Không phải tôi không hoan nghênh chuyện này. Nhưng mà\dots\ mọi thứ đến quá nhanh. Quá bất ngờ. Đến mức mà không một ai trong gia đình chúng tôi đều ngờ tới.

Mới ngày hôm qua, Aloe-chan vẫn còn lang bạt đâu đó, cố gắng vực dậy tinh thần sau một chuỗi sự kiện không hay. Vậy mà hôm nay, em ấy đã có một căn phòng rộng rãi, tiện nghi hơn cả phòng của tôi.

Một phần trong tôi cảm thấy nhẹ nhõm: dù gì, Aloe-chan cũng còn quá trẻ để đối mặt với mọi thứ một mình. Nếu nhà Hikawa có thể giúp em ấy có một khởi đầu mới, thì đó là điều đáng mừng.

Nhưng đồng thời, tôi cũng không thể không đặt dấu hỏi về mục đích của Tanigo-san khi lập ra nơi này.

Ban đầu, tôi cứ nghĩ đó đơn thuần là một sự tử tế theo kiểu ông chú tốt bụng mà YAGOO luôn thể hiện. Nhưng càng nghĩ, tôi càng thấy có điều gì đó sâu xa hơn.

Là để kết nối lại các thành viên đã rời đi với \hll?

Hay là để tạo ra một nơi mà những ai từng gắn bó với \hll\ có thể tìm về khi họ mất phương hướng?

Nếu đúng như thế, thì\dots\ Aloe có lẽ không phải là người duy nhất sẽ đặt chân đến đây.

Tôi khẽ thở dài, gạt mớ suy nghĩ ấy sang một bên. Dù lý do là gì đi nữa, trước mắt tôi vẫn có một việc quan trọng hơn - đó là đảm bảo em ấy có thể ổn định cuộc sống ở đây lâu dài.

Một cái Tết nữa sắp tới rồi. Lần này, tôi không chỉ đón những ngày Tết Nguyên Đán với gia đình, mà còn đón chung với một thành viên mới. Không biết liệu em ấy có thấy xa lạ không nhỉ?

\il

\pov{Mano Aloe}

Bên trong căn phòng mới, ta lặng lẽ nhìn ra cửa sổ. Bầu trời mùa đông xám xịt, nhưng ánh đèn từ các tòa nhà gần đó vẫn nhấp nháy như một lời nhắc rằng thế giới vẫn đang chuyển động, bất kể ta đang ở phương trời nào đi nữa.

Có lẽ, đây là lần đầu tiên sau nhiều tháng, ta thực sự có một nơi để gọi là nhà. Một nơi không chỉ để lưu trú, mà là nơi để thực sự cảm nhận được bầu không khí của gia đình.

Cuộc gọi bất ngờ của YAGOO ngày hôm đó quả thực đã thay đổi cuộc đời ta sang một trang mới.

Một lời đề nghị ta chưa từng nghĩ đến, một nơi ở mới, một cơ hội để bắt đầu lại từ đầu. Không phải là một quán trọ tạm bợ, cũng không phải là một góc nhỏ trong nhà bạn bè. Mà là một nơi với những con người mà, dù xa lạ, lại khiến ta cảm thấy được chào đón.

Tại sao ông ta lại làm thế? Đó có phải chỉ là lòng tốt đơn thuần, là sự đền bù của tôi sau những ngày tháng vất vả đó?

Không, không thể thế được. Ông ta không phải kiểu người sẽ đưa ra quyết định lớn như vậy mà không có lý do.

Ông ấy đang cố kết nối mọi người lại với nhau?
Hay là đang chuẩn bị cho điều gì đó trong tương lai?

Thực sự, đến thời điểm này, ta cũng không biết.

Chỉ biết rằng, nếu đây là cơ hội để tôi có thể hội ngộ với các thành viên trong Gen 5 - những đồng đội cũ đã từng hỉ nộ ái ố với ta - thì đây sẽ là một cơ hội không thể nào từ chối.

Nhìn căn phòng giờ đã trở nên quen thuộc hơn, ta khẽ mỉm cười. Năm nay sẽ khác. Ta không còn đơn độc nữa.

Những ngày gần đây, nhân gian đang bắt đầu hối hả hơn. Có vẻ như họ đang chuẩn bị một kỳ nghỉ dài nào đó\dots

Tết, đúng không nhỉ?

Nếu vậy, đây có thể sẽ là năm đầu tiên ta đón một cái Tết như vậy. Có lẽ ta cũng nên thử đón ngày này với tư cách là một người trong gia đình\dots

Sẽ rất là vui đây!

\end{document}
